\documentclass[UTF8,12pt,a4paper,openany]{ctexbook}
\usepackage[margin=2.2cm]{geometry}
\usepackage{amsmath,amssymb}
\usepackage{enumitem}
\usepackage[most]{tcolorbox}
\usepackage{xcolor}
\usepackage{hyperref}
\hypersetup{colorlinks=true,linkcolor=blue}
\newcommand{\TODO}[1]{\textcolor{red}{[需校对：#1]}}
\tcbset{colback=gray!3,colframe=black!45,arc=2mm,boxrule=0.5pt}
\title{结构力学十四讲作业、解析与答案\\\large OCR 转写 LaTeX 重排版}
\author{根据本地 PDF 资料转写整理}
\date{\today}
\begin{document}
\maketitle
\frontmatter
\chapter*{使用说明}
本版不再采用整页扫描图嵌入，而是将题干、答案和解析转为可编辑 LaTeX 文本。原始资料中部分题面为扫描图片，部分讲解为手写内容，OCR 对结构图、复杂公式和手写字识别不稳定；文中以 \TODO{...} 标出需要人工对照原 PDF 校对的位置。
\tableofcontents
\mainmatter
\chapter{第1讲：几何组成分析}
\begin{tcolorbox}[title=解析总览]
先判定体系自由度，再结合二元体规则、两刚片规则、三刚片规则和瞬变判据作结论。
\end{tcolorbox}
\section{答案：结构力学基础与技巧班第1次作业答案（几何组成分析）}
\begin{enumerate}[label=\textbf{\arabic*.},leftmargin=2em]
\item \TODO{OCR 自动转写，结构图/公式需对照原 PDF 校对。}\par 【第 1 页，来源：pdf-text】
\item \TODO{OCR 自动转写，结构图/公式需对照原 PDF 校对。}\par 1-1、D\par 解析：采用两刚片规则分析，由于4根链杆平行且不等长，因此为瞬变体系，有\par 两个多余约束。
\item \TODO{OCR 自动转写，结构图/公式需对照原 PDF 校对。}\par 1-2、B\par 解析：原体系的计算自由度\par 3 3 1 2\par 4\par 3\par W =\textbackslash{}\textbackslash{}times -\textbackslash{}\textbackslash{}times -\par =\par ，故需要再补充三个约束。当\par 在A处加入固定端时，刚片ABC先并入大地，随后刚片DE与大地由链杆BD和铰\par E连接，满足两刚片规则可并入大地，同理FG与大地同样满足两刚片规则，此时\par 原体系变为无多余约束的几何不变体系。
\item \TODO{OCR 自动转写，结构图/公式需对照原 PDF 校对。}\par 1-3、C\par 解析：刚片ABC与大地由三根不共线且不交于一点的链杆A、BD和CE，满足两\par 刚片规则，原体系为无多余约束的几何不变体系。
\item \TODO{OCR 自动转写，结构图/公式需对照原 PDF 校对。}\par 1-4、几何常变；1\par 解析：采用两刚片规则分析，刚片ABCD和有一个多余约束的刚片EFGH仅由两\par 根链杆CE和DF连接，不满足两刚片规则，原体系为有一个多余约束的几何常变\par 体系。\par \par 【第 2 页，来源：pdf-text】
\item \TODO{OCR 自动转写，结构图/公式需对照原 PDF 校对。}\par 1-5、-4\par 解析：刚片如图所示，计算自由度\par 5 3\par 6 2 1 3\par 4\par 4\par W =\textbackslash{}\textbackslash{}times -\textbackslash{}\textbackslash{}times\par -\textbackslash{}\textbackslash{}times -\par =-。
\item \TODO{OCR 自动转写，结构图/公式需对照原 PDF 校对。}\par 1-6、×\par 解析：由大地依次增加二元体AB、ACB、ADC、DEC、DFE、EGH、EIG、FJI、\par FKJ、KLF、KML，原体系为无多余约束的几何不变体系。
\item \TODO{OCR 自动转写，结构图/公式需对照原 PDF 校对。}\par 1-7、解析：刚片AD、BE可直接并入大地，链杆DE为一个多余约束。随后刚片\par EFC与大地由固定端C和铰E连接，满足两刚片规则且有两个多余约束，故原体\par 系为有3个多余约束的几何不变体系。\par \par 【第 3 页，来源：pdf-text】
\item \TODO{OCR 自动转写，结构图/公式需对照原 PDF 校对。}\par 1-8、解析：采用三刚片规则分析，各刚片如图所示，刚片ABC与刚片CDE交于\par 铰C，刚片ABC与刚片GF由链杆BF、AG连接交于铰\par 2\par O ，刚片CDE与刚片FG由\par 链杆DF、GE连接交于铰\par 1\par O ，三铰不共线满足三刚片规则组成大刚片，后与大地\par 由三根不共线且不交于一点的链杆支座连接，满足两刚片规则，原体系为无多余\par 约束的几何不变体系。
\item \TODO{OCR 自动转写，结构图/公式需对照原 PDF 校对。}\par 1-9、解析：采用三刚片规则分析，各刚片如图所示，刚片ABC与刚片DEF由链\par 杆BD、CF连接交于无穷远铰\par 1\par O ，刚片ABC与大地刚片由链杆A、CG连接交于铰\par A，刚片DEF与大地刚片由链杆FG、E连接交于铰E，三铰共线不满足三刚片规\par 则，原体系为有一个多余约束的几何瞬变体系。
\item \TODO{OCR 自动转写，结构图/公式需对照原 PDF 校对。}\par 1-10、解析：计算自由度\par 12 2\par 25\par 1\par W =\par \textbackslash{}\textbackslash{}times -\par =-。采用两刚片规则分析，刚片ABCD\par 与刚片DEFB由铰B和铰D共四个约束连接，满足两刚片规则且有一个多余约束，\par 原体系为有一个多余约束的几何不变体系。\par \par 【第 4 页，来源：pdf-text】
\item \TODO{OCR 自动转写，结构图/公式需对照原 PDF 校对。}\par 1-11、解析：先去掉二元体AB、CEG，随后有两个多余约束的刚片BCD与大地\par 由三根不共线且不交于一点的链杆D、CF和CG连接，满足两刚片规则，原体系\par 为有两个多余约束的几何不变体系，两次超静定。
\item \TODO{OCR 自动转写，结构图/公式需对照原 PDF 校对。}\par 1-12、解析：刚片ABC与大地由三根不共线且不交于一点的链杆支座A和链杆CD\par 连接，满足两刚片规则可并入大地，随后刚片FBE与大地同样由铰B和链杆F连接\par 满足两刚片规则，杆GE为多余约束。原体系为有一个多余约束的几何不变体系。
\item \TODO{OCR 自动转写，结构图/公式需对照原 PDF 校对。}\par 1-13、解析：可将刚片3456等价为链杆34，如图所示。随后采用两刚片规则分析，\par 刚片14与大地由三根不共线且不交于一点的链杆连接，满足两刚片规则，原体系\par 为无多余约束的几何不变体系。\par \par 【第 5 页，来源：pdf-text】
\item \TODO{OCR 自动转写，结构图/公式需对照原 PDF 校对。}\par 1-14、解析：先去掉二元体ABC、AJ、DEF和GFH，随后刚片JCDI与大地由三根\par 不共线且不交于一点的链杆支座连接，满足两刚片规则，原体系为无多余约束的\par 几何不变体系。
\item \TODO{OCR 自动转写，结构图/公式需对照原 PDF 校对。}\par 1-15、解析：计算自由度\par 7 2 13 1\par 1\par W =\par \textbackslash{}\textbackslash{}times -\par \textbackslash{}\textbackslash{}times =。采用两刚片规则分析，刚片ABC\par 与刚片DEF仅由链杆CD、BE连接，不满足两刚片规则，原体系为无多余约束的\par 几何常变体系。
\item \TODO{OCR 自动转写，结构图/公式需对照原 PDF 校对。}\par 1-16、解析：先去掉二元体ABC，随后采用三刚片规则分析，各刚片如图所示，\par 刚片ADE与大地刚片由链杆D、EH连接交于无穷远铰\par 1\par O ，刚片CFG与大地刚片\par 由链杆F、GH连接交于铰\par 2\par O ，刚片ADE与刚片CFG由链杆AC、EG连接交于无穷\par 远铰\par 3\par O ，三铰共线不满足三刚片规则，原体系为有一个多余约束的几何瞬变体\par 系。\par \par 【第 6 页，来源：pdf-text】
\item \TODO{OCR 自动转写，结构图/公式需对照原 PDF 校对。}\par 1-17、解析：计算自由度\par 4 3\par 3 2\par 7 1\par 1\par W =\par \textbackslash{}\textbackslash{}times -\textbackslash{}\textbackslash{}times\par -\textbackslash{}\textbackslash{}times =-。刚片AB可先并入大地，随\par 后采用两刚片规则分析，杆CD与大地由三根不共线且不交于一点的链杆BC和支\par 座D连接，满足两刚片规则，CE杆为多余约束，故原体系为有一个多余约束的几\par 何不变体系。
\end{enumerate}
\section{作业题目：结构力学基础与技巧班第1次作业（几何组成分析）}
\begin{enumerate}[label=\textbf{\arabic*.},leftmargin=2em]
\item \TODO{OCR 自动转写，结构图/公式需对照原 PDF 校对。}\par 【第 1 页，来源：ocr】\par 权力矢础焉第/次作业一 一 IETem ena      888\par 全7 有人和第(次作业 全人区全合作91680888\par 几何组成分析作业及答案
\item \TODO{OCR 自动转写，结构图/公式需对照原 PDF 校对。}\par 1-1 1- 56 所示体系的几何组成为  )。(浙江大学 2008)\par AL 几何瞬变, 无多余约束             B. 几何不变\par C. 几何常变                   D. 几何瞬变，有多余约束
\item \TODO{OCR 自动转写，结构图/公式需对照原 PDF 校对。}\par 1-2 (EA 1 - 57 所示体系成为无多余约束的几何不变体系，则需在A端加入 CD.\par (中南大学 2011，合肥工业大学 2009)\par AL FE BESTE   B. 固定支座     C. RE    D. 定向支座
\item \TODO{OCR 自动转写，结构图/公式需对照原 PDF 校对。}\par (1-3 图1-58所示体系的几何构造性质是 〈  )。(中南大学 2009)\par A. 常变体系                   B，盟变体系\par C. 无多余联系的几何不变体系       D. 有多余联系的几何不变体系\par 人四 | |\par 图1- 56               图1-57              图1-58
\item \TODO{OCR 自动转写，结构图/公式需对照原 PDF 校对。}\par 1-4 图1-59所示体系内部是体系, 它有\_ 个多余约束。 (广西大学\par 2010)
\item \TODO{OCR 自动转写，结构图/公式需对照原 PDF 校对。}\par 1-5 图1-60所示体系的计算自由度 要一       。(了哈尔滨工业大学 2011)
\item \TODO{OCR 自动转写，结构图/公式需对照原 PDF 校对。}\par 1-6 图1-61 所示体系几何不变上且有多余联系。(  ) (中南大学 2012)\par 、、, 44. 7/\par ci\textbackslash{}@\par 第18页,大359页                                 Gaede FRED\par \par 【第 2 页，来源：ocr】\par ROR OBIGSIBSS\par fl 1-59                FA 1-60               图1-61     .     -
\item \TODO{OCR 自动转写，结构图/公式需对照原 PDF 校对。}\par 1-7 XH 1-62 所示体系进行几何组成分析。(哈尔滨工业大学 2009)
\item \TODO{OCR 自动转写，结构图/公式需对照原 PDF 校对。}\par 1-8 对图1-63 所示体系作几何组成分析。(武汉科技大学 2009)
\item \TODO{OCR 自动转写，结构图/公式需对照原 PDF 校对。}\par 1-9 对图1-64所示杆件体系进行几何构造分析。(武汉理工大学 2008)\par D                E             F                                                                                      .\par A            |           |                                                     SS  “A\par 图1-62                图1-63        图1-64
\item \TODO{OCR 自动转写，结构图/公式需对照原 PDF 校对。}\par 1-10 计算图1-65 所示体系的计算自由度，  试分析其体系的几何组成。 (华南理工大-\par 学 2012)
\item \TODO{OCR 自动转写，结构图/公式需对照原 PDF 校对。}\par 1-11 对图1- 66 所示平面体系进行几何组成分析，  并指出结构超静定次数。  (青岛理\par 工大学 2012)
\item \TODO{OCR 自动转写，结构图/公式需对照原 PDF 校对。}\par 1-65 1-66
\item \TODO{OCR 自动转写，结构图/公式需对照原 PDF 校对。}\par 1-12 分析图1- 67 所示体系的几何稳定性，写出分析过程。(东南大学2008) —
\item \TODO{OCR 自动转写，结构图/公式需对照原 PDF 校对。}\par 1-13 试对图 1-68 所示体系作几何组成分析。(山东建筑大学 2011)\par ]                     |           1         2      4                         3\par ie                                 5            6
\item \TODO{OCR 自动转写，结构图/公式需对照原 PDF 校对。}\par 1-67.                                       图1-68\par vat Va.,\par O74 \textbackslash{}@\par 第19页,大359页                                 aa aie FERED\par Bi-:RE ASA\par \par 【第 3 页，来源：ocr】\par EEA Rains)\par ROGUE WBIGHTSSS
\item \TODO{OCR 自动转写，结构图/公式需对照原 PDF 校对。}\par 1-14 试对图 1-69 所示体系作几何组成分析。(山东建筑大学 2008)\par «1-15 计算图1-70 的计算自由度，并作几何构造分析。(中国地震局工程力学研究所\par 2010)\par 图1-69                               ，图1-70
\item \TODO{OCR 自动转写，结构图/公式需对照原 PDF 校对。}\par 1-16 对图 1- 71 所示体系进行几何构造分析。(西南交通大学 2008)
\item \TODO{OCR 自动转写，结构图/公式需对照原 PDF 校对。}\par 1-17 计算图 1-72 所示体系的计算自由度，并进行几何构成分析。 (北京工业大学\par 2010)                                                                        、\par A              ..B\par | [4     aa\par , 1-71 1-72                            .\par wla.s\par RONG\par SON)\par 第20页, F359 Hk                                 faa LED\par Bi:ReE Asia
\end{enumerate}
\chapter{第2讲：理论力学回顾}
\begin{tcolorbox}[title=解析总览]
先画隔离体受力图，再分别列水平力、竖向力和力矩平衡方程。
\end{tcolorbox}
\section{答案：结构力学基础与技巧班第2次作业答案（理论力学回顾）}
\begin{enumerate}[label=\textbf{\arabic*.},leftmargin=2em]
\item \TODO{OCR 自动转写，结构图/公式需对照原 PDF 校对。}\par 【第 1 页，来源：pdf-text】\par (a)解：取整体隔离体竖向力平衡，可得B处竖向支反力\par [符号][符号]\par 0\par 0\par 1\par 1\par 2\par 2\par By\par F\par l\par q\par q l\par =\par \textbackslash{}\textbackslash{}times \textbackslash{}\textbackslash{}times\par =\par [符号]\par 对B点取矩得B处支座反力矩\par 2\par 0\par 0\par 1\par 1\par 1\par 2\par 3\par 6\par B\par M\par l\par q\par l\par q l\par =\par \textbackslash{}\textbackslash{}times \textbackslash{}\textbackslash{}times\par \textbackslash{}\textbackslash{}times\par =\par (顺时针)\par (b)解：取整体隔离体竖向力平衡，可得A处竖向支反力\par [符号][符号]\par 15 1\par 30\par 45\par Ay\par F\par KN\par =\par \textbackslash{}\textbackslash{}times +\par =\par [符号]\par 对A点取矩得A处支座反力矩\par 15 1 2.5\par 30 3\par 127.5\par A\par M\par KN m\par =\par \textbackslash{}\textbackslash{}times \textbackslash{}\textbackslash{}times\par +\par \textbackslash{}\textbackslash{}times =\par \textbackslash{}\textbackslash{}cdot\par (逆时针)\par (c)解：取整体隔离体对C点取矩，由\par 3 30 40 1.5\par 0\par B\par F \textbackslash{}\textbackslash{}times +\par \textbackslash{}\textbackslash{}times\par =\par ，可得B点竖向支\par 反力\par 10\par B\par F\par KN\par =\par ，竖向力平衡得\par 30\par C\par F\par KN\par =\par (d)解：取整体隔离体对A点取矩，由\par 10 4 0.2 10 5 0\par B\par F \textbackslash{}\textbackslash{}times\par \textbackslash{}\textbackslash{}times\par \textbackslash{}\textbackslash{}times =，可得B点竖向\par 支反力\par 1.4\par B\par F\par KN\par =\par ，竖向力平衡得\par 0.6\par A\par F\par KN\par =\par (e)解：取整体隔离体对A点取矩，由\par 3 6 20 2 20\par 0\par B\par F \textbackslash{}\textbackslash{}times --\par \textbackslash{}\textbackslash{}times -\par =\par ，可得B点竖向\par 支反力\par 22\par B\par F\par KN\par =\par ，竖向力平衡得\par 2\par A\par F\par KN\par =\par \par 【第 2 页，来源：pdf-text】\par (f)解：取整体隔离体对A点取矩，由\par 4\par 4.5\par 0\par C\par F\par a q a\par a\par \textbackslash{}\textbackslash{}times\par -\textbackslash{}\textbackslash{}times \textbackslash{}\textbackslash{}times\par =\par ，可得B点竖向支\par 反力\par 1.125\par C\par F\par qa\par =\par ，竖向力平衡得\par 0.125\par A\par F\par qa\par =\par (g)解：取整体隔离体竖向力平衡，可得A处竖向支反力\par [符号][符号]\par 5 1\par 5\par Ay\par F\par KN\par =\par \textbackslash{}\textbackslash{}times =\par [符号]\par 对A点取矩得A处支座反力矩\par 5 1 2\par 10\par A\par M\par KN m\par =\textbackslash{}\textbackslash{}times \textbackslash{}\textbackslash{}times\par =\par \textbackslash{}\textbackslash{}cdot\par (逆时针)\par (h)解：取整体隔离体竖向力平衡，可得A处竖向支反力\par [符号][符号]\par 15\par Ay\par F\par KN\par =\par [符号]\par 对A点取矩得A处支座反力矩\par 10 15 1\par 25\par A\par M\par KN m\par =\par +\par \textbackslash{}\textbackslash{}times =\par \textbackslash{}\textbackslash{}cdot\par (逆时针)\par (i)解：取整体隔离体对A点取矩，由\par 5\par 3\par 2.5\par 0\par B\par F\par a q\par a\par a\par \textbackslash{}\textbackslash{}times\par -\textbackslash{}\textbackslash{}times\par \textbackslash{}\textbackslash{}times\par =，可得B点竖向支\par 反力\par 1.5\par B\par F\par qa\par =\par ，竖向力平衡得\par 1.5\par A\par F\par qa\par =\par \par 【第 3 页，来源：pdf-text】\par (j)解：取整体隔离体对A点取矩，由\par 3\par 2\par 0\par B\par e\par e\par F\par a M\par M\par \textbackslash{}\textbackslash{}times\par +\par =\par ，可得B点竖向支反\par 力\par 3\par e\par B\par M\par F\par a\par =\par ，竖向力平衡得\par 3\par e\par A\par M\par F\par a\par =\par (k)解：取整体隔离体对A点取矩，可得\par 2\par 2\par By\par Bx\par F\par a\par F\par a\par q\par a a\par \textbackslash{}\textbackslash{}times\par +\par \textbackslash{}\textbackslash{}times =\textbackslash{}\textbackslash{}times\par \textbackslash{}\textbackslash{}times\par ；取BC部\par 分隔离体对C点取矩可得\par By\par Bx\par F\par a\par F\par a\par \textbackslash{}\textbackslash{}times =\par \textbackslash{}\textbackslash{}times ，联立解得\par 2\par 3\par By\par F\par qa\par =\par 2\par 3\par Bx\par F\par qa\par =\par 。再由\par 整体竖向力和水平力平衡，可得A处支座反力\par 2\par 3\par Ay\par F\par qa\par =\par 4\par 3\par Ax\par F\par qa\par =\par (l)解：取整体隔离体对A点取矩，可得\par 10\par 2 5 7.5\par 7.5\par By\par By\par F\par F\par KN\par \textbackslash{}\textbackslash{}times\par =\textbackslash{}\textbackslash{}times \textbackslash{}\textbackslash{}times\par [符号]\par =\par ；取\par BC部分隔离体对C点取矩可得\par 25\par 5\par 6\par 2 5 2.5\par 12\par By\par Bx\par Bx\par F\par F\par F\par KN\par \textbackslash{}\textbackslash{}times\par =\par \textbackslash{}\textbackslash{}times +\textbackslash{}\textbackslash{}times \textbackslash{}\textbackslash{}times\par [符号]\par =\par 。再由整\par 体竖向力和水平力平衡，可得A处支座反力\par 2.5\par Ay\par F\par KN\par =\par 25\par 12\par Ax\par F\par KN\par =\par \par 【第 4 页，来源：pdf-text】\par (m)解：取整体隔离体对A点取矩，可得\par 2\par 1\par 2\par 2\par Dy\par Dy\par qh\par F\par l\par q\par h\par h\par F\par l\par \textbackslash{}\textbackslash{}times =\par \textbackslash{}\textbackslash{}times\par \textbackslash{}\textbackslash{}times\par [符号]\par =\par ，再由\par 整体竖向力和水平力平衡，可得A处支座反力\par 2\par 2\par Ay\par qh\par F\par l\par =\par Dx\par F\par qh\par =
\end{enumerate}
\section{作业题目：结构力学基础与技巧班第2次作业（理论力学回顾）}
\begin{enumerate}[label=\textbf{\arabic*.},leftmargin=2em]
\item \TODO{OCR 自动转写，结构图/公式需对照原 PDF 校对。}\par 【第 1 页，来源：ocr】\par \$ hey ih 1644 aa 1 BH\par 30 kN\par i                         15 kN/m\par *                          me       a                 C    lm\par |一一一一\par (a)                              (b)\par 4kN+m\par 30kN-m             40 kN                0.2 kNym ”从\par 10m\par (c)                              (d)\par 6kN-m             20 kKN-m 20 kN                           q\par im                          Sa\par 3m                  | .         sa\par (e)                              (f)\par YA A                              本\par |          2.5m                         lm\par 2m\par G)                              (h)\par )           q                        iM,                   2M,\par A                        RB      P       om                )\par Cc           D                    C\par ‘a        3a             ,         2 a                      ?\par a “aa                     3a\par Y                                、\par NS\par 255, SSS\par O75 \textbackslash{}O\par BA Sete FIRED\par Bi-RE ASR\par \par 【第 2 页，来源：ocr】\par ，  BE eae erties\par 。      DY IRSASENIGILSSII DSOSS\par q                      上\par 2           ry\par Bf          |  5m \_ sh  3m |\par B      CG\par A    D ,\par (m)\par i 5a LED\par Bit-:RE ize
\end{enumerate}
\chapter{第3讲：材料力学回顾}
\begin{tcolorbox}[title=解析总览]
按内力、应力、变形和强度/刚度条件组织解答，注意单位与符号。
\end{tcolorbox}
\section{答案：结构力学基础与技巧班第3次作业答案（材料力学回顾）}
\begin{enumerate}[label=\textbf{\arabic*.},leftmargin=2em]
\item \TODO{OCR 自动转写，结构图/公式需对照原 PDF 校对。}\par 【第 1 页，来源：pdf-text】\par 一、\par (a)解：取整体隔离体竖向力平衡，可得B处竖向支反力\par [符号][符号]\par 0\par 0\par 1\par 1\par 2\par 2\par By\par F\par l\par q\par q l\par =\par \textbackslash{}\textbackslash{}times \textbackslash{}\textbackslash{}times\par =\par [符号]\par 对B点取矩得B处支座反力矩\par 2\par 0\par 0\par 1\par 1\par 1\par 2\par 3\par 6\par B\par M\par l\par q\par l\par q l\par =-\par \textbackslash{}\textbackslash{}times \textbackslash{}\textbackslash{}times\par \textbackslash{}\textbackslash{}times\par =-\par (上侧受拉)\par 画出原结构弯矩图和剪力图如图(a)(b)所示。\par (b )解：取整体隔离体竖向力平衡，可得A 处竖向支反力\par 15 1 30\par Ay\par F\par =\par \textbackslash{}\textbackslash{}times +\par =\par [符号][符号]\par 45KN [符号]\par ，对A 点取矩得A 处支座反力矩\par 15 1 2.5\par 30 3\par A\par M\par =-\par \textbackslash{}\textbackslash{}times \textbackslash{}\textbackslash{}times\par \textbackslash{}\textbackslash{}times =\par 127.5KN m\par \textbackslash{}\textbackslash{}cdot\par (上侧受拉)。弯矩图如图(a)所示。\par (c)解：取整体隔离体对C点取矩，由\par 3 30 40 1.5 0\par B\par F \textbackslash{}\textbackslash{}times +\par \textbackslash{}\textbackslash{}times\par =\par ，可得B点竖向支\par 反力\par 10\par B\par F\par KN\par =\par ，竖向力平衡得\par 30\par C\par F\par KN\par =\par 。画出弯矩图如图(b)(c)所示。\par \par 【第 2 页，来源：pdf-text】\par (d)解：取整体隔离体对A点取矩，由\par 10 4 0.2 10 5 0\par B\par F \textbackslash{}\textbackslash{}times\par \textbackslash{}\textbackslash{}times\par \textbackslash{}\textbackslash{}times =，可得B点竖向\par 支反力\par 1.4\par B\par F\par KN\par =\par ，竖向力平衡得\par 0.6\par A\par F\par KN\par =\par 。C处作用了集中力偶，两侧截面\par 的弯矩有突变，由C左侧截开取AC隔离体，可得\par 0.6 8 0.2 8 4\par L\par C\par M =\par \textbackslash{}\textbackslash{}times -\par \textbackslash{}\textbackslash{}times \textbackslash{}\textbackslash{}times =1.6KN m\par \textbackslash{}\textbackslash{}cdot\par (上侧受拉)，由C结点力矩平衡得\par 2.4\par R\par C\par M\par KN m\par =\par \textbackslash{}\textbackslash{}cdot\par (下侧受拉)。画出弯矩图和剪力图如图(b)(c)所示。\par (e)解：取整体隔离体对A点取矩，由\par 3 6 20 2 20\par 0\par B\par F \textbackslash{}\textbackslash{}times --\par \textbackslash{}\textbackslash{}times -\par =\par ，可得B点竖向\par 支反力\par 22\par B\par F\par KN\par =\par ，竖向力平衡得\par 2\par A\par F\par KN\par =\par 。画出弯矩图和剪力图如图(b)(c)\par 所示。\par \par 【第 3 页，来源：pdf-text】\par (f)解：取整体隔离体对A点取矩，由\par 4\par 4.5\par 0\par C\par F\par a q a\par a\par \textbackslash{}\textbackslash{}times\par -\textbackslash{}\textbackslash{}times \textbackslash{}\textbackslash{}times\par =，可得B点竖向支\par 反力\par 1.125\par C\par F\par qa\par =\par ，竖向力平衡得\par 0.125\par A\par F\par qa\par =\par 。画出弯矩图和剪力图如图(b)\par (c)所示。\par 二、\par (a)解：取整体隔离体竖向力平衡，可得A处竖向支反力\par [符号][符号]\par 5 1\par 5\par Ay\par F\par KN\par =\par \textbackslash{}\textbackslash{}times =\par [符号]\par 对A点取矩得A处支座反力矩\par 5 1 2\par 10\par A\par M\par KN m\par =-\textbackslash{}\textbackslash{}times \textbackslash{}\textbackslash{}times =-\par \textbackslash{}\textbackslash{}cdot\par (上侧受拉)\par 画出弯矩图和剪力图如图(a)(b)所示。\par (b)解：取整体隔离体竖向力平衡，可得A处竖向支反力\par [符号][符号]\par 15\par Ay\par F\par KN\par =\par [符号]\par 对A点取矩得A处支座反力矩\par 10 15 1\par 25\par A\par M\par KN m\par =-\par \textbackslash{}\textbackslash{}times =-\par \textbackslash{}\textbackslash{}cdot\par (上侧受拉)\par 画出弯矩图和剪力图如图(a)(b)所示。\par \par 【第 4 页，来源：pdf-text】\par (c)解：取整体隔离体对A点取矩，由\par 5\par 3\par 2.5\par 0\par B\par F\par a q\par a\par a\par \textbackslash{}\textbackslash{}times\par -\textbackslash{}\textbackslash{}times\par \textbackslash{}\textbackslash{}times\par =，可得B点竖向\par 支反力\par 1.5\par B\par F\par qa\par =\par ，竖向力平衡得\par 1.5\par A\par F\par qa\par =\par 。画出弯矩图和剪力图如图(b)(\par c)所示。\par (d)解：取整体隔离体对A点取矩，由\par 3\par 2\par 0\par B\par e\par e\par F\par a M\par M\par \textbackslash{}\textbackslash{}times\par +\par =\par ，可得B点竖向支\par 反力\par 3\par e\par B\par M\par F\par a\par =\par ，竖向力平衡得\par 3\par e\par A\par M\par F\par a\par =\par 。画出弯矩图和剪力图如图(b)(c)所\par 示。\par \par 【第 5 页，来源：pdf-text】\par (e)解：取整体隔离体对A点取矩，由\par 4 4 4 2 4 20\par B\par F \textbackslash{}\textbackslash{}times =\textbackslash{}\textbackslash{}times \textbackslash{}\textbackslash{}times ++\par 可得B处支座反力\par 14\par B\par F\par KN\par =\par 。再由整体竖向力平衡，可得A处支座反力\par 2\par A\par F\par KN\par =\par 。画出弯矩图和\par 剪力图如图(b)(c)所示。\par (f)解：取整体隔离体对A点取矩，由\par 11\par 4\par 2\par 3\par 8\par B\par F\par a\par F\par a\par F\par a\par F\par a\par \textbackslash{}\textbackslash{}times\par +\par \textbackslash{}\textbackslash{}times\par =\par \textbackslash{}\textbackslash{}times\par +\par \textbackslash{}\textbackslash{}times\par ，可得B\par 处支座反力\par 5\par 16\par B\par F\par F\par =\par ，再由整体竖向力平衡，可得A处支座反力\par 5\par 16\par A\par F\par F\par =\par 。画\par 出弯矩图和剪力图如图(b)(c)所示。\par \par 【第 6 页，来源：pdf-text】\par (g)解：取整体隔离体对A点取矩，由\par 2\par 2 1 0.5\par B\par F \textbackslash{}\textbackslash{}times =\textbackslash{}\textbackslash{}times \textbackslash{}\textbackslash{}times\par 可得B结点的竖向反力\par 1\par 2\par B\par F\par KN\par =\par ，再由整体竖向力平衡，可得A处支座反力\par 3\par 2\par A\par F =\par 。画出弯矩图和剪\par 力图如图(b)(c)所示。\par (h)解：取整体隔离体对A点取矩，由\par 6 10 6 3 250 2 250 4\par B\par F \textbackslash{}\textbackslash{}times =\par \textbackslash{}\textbackslash{}times \textbackslash{}\textbackslash{}times +\par \textbackslash{}\textbackslash{}times +\par \textbackslash{}\textbackslash{}times 可得B\par 处支座的竖向反力\par 280\par B\par F\par KN\par =\par ，再由整体竖向力平衡，可得A处支座反力\par 280\par A\par F\par KN\par =\par 。画出弯矩图和剪力图如图(b)(c)所示。\par 三、\par \par 【第 7 页，来源：pdf-text】\par (a)解：取CD部分隔离体并对B点取矩，由\par C\par F\par a\par qa a\par \textbackslash{}\textbackslash{}times =\par \textbackslash{}\textbackslash{}times 可得C结点的竖向反\par 力\par C\par F\par qa\par =\par ，再取AC部分，可得A处固定端反力矩\par 2\par 2\par 1\par 1\par 2\par 2\par A\par M\par qa a\par qa\par qa\par =\par \textbackslash{}\textbackslash{}times\par =\par (下\par 侧受拉)。画出弯矩图和剪力图如图(b)(c)所示。\par (b)解：已知AC为二力杆，没有剪力和弯矩，易画出弯矩图和剪力图如图(a)\par (b)所示。
\end{enumerate}
\section{作业题目：结构力学基础与技巧班第3次作业（材料力学回顾）}
\begin{enumerate}[label=\textbf{\arabic*.},leftmargin=2em]
\item \TODO{OCR 自动转写，结构图/公式需对照原 PDF 校对。}\par 【第 1 页，来源：ocr】\par 小鹿土木考研\par 第3? 次作业一一材料力学回碳\par 材料力学回顾作业及答案\par 一、男前力图和弯矩图\par 30kN\par TH                   15 kN/m |\par 4                    .     4            C   1 m   2\par -一一一一一        am\par (a)                         (b)\par 4kN-m\par 30kN-m           40 kN              0.2kN/m\_ 7)\par C        D         A\par A                                                B\par lm |        1.5m                10m\par 4m\par (c)                         (d)\par 6kKN-m KN-m 90 kN                             q\par Cs             C       BY)  人      —— fit     iB\par Im                       Sa        2\par |         3m                 .           加  aan\par (e)                         (f)\par ry 1 1 。 7/\par N  A\par anes\par 小鹿土木考研第 43/359 页                            ee 四\par Ai VO\par BA ER CRED\par BSit-Re AiR\par \par 【第 2 页，来源：ocr】\par 一 (>  可wo    \textbackslash{} Qs LEM fy 4a tcl 5\par DEAT ODL) SER Op anr ame\par 二、 试利用荷载集度 .剪力和弯矩间的微分关系作下列各梁的前力图和弯和矩图。\par 5 kN/m               15 kN ,10kN-m\par C     lm    7     4\}            | Ss\par 2m\par (a)                           (b)\par q                     M,                 OM,\par Cc          D            BS                ¢\par ‘a     3a         °      2 CQ              2\par Sa                             3a\par (c)                           (d)\par A                       B      ()    c    Dit    z   -\par |\}\_\_\_\_4m\_\_.                  4a\par (e)                           (f)\par att                 10 KN/m J 250 KN 4250 KN\par a. =      i\par 5                A\par 6m\par (g)                           (h)\par = Gah HEE Ba SLAs ee A BY J BS\par yt          oa            qa   M=qa.\par ACP BD, «ARP BE >)\par Lelepa     [La | 4 lia |\par (a)                      (b)\par . 1 1 。/\par YU\par 土木考研第 44/359 页                           Ses ee\par fA ee FRED\par \par 【第 3 页，来源：ocr】\par FEARS OE     —\par 去六合信1981387888\par (0)\par Warts Aathged a-ha,\par 下         f,\par 了       J          at\par 1       A                                '\par (a-')\par 2在向的\par ENMezo re a 4.4  中 7  4, ,  \textbackslash{}\par 3    一           本   一   zx\par sh          oe\par (a-2)                        6-5)\par Jo\par 出           is\par A j\_2 Gy\par pal\par as 9 Ha tas fa H\%-\par HAL 2对- y aon ihe\par Is   J    出 tf0, Qzsxitsoz 外Mo\par AP ?             乡      \_\par 720, M +Jox (7-1) + ISX lx (\$4) 20\par Rs 4                M= BS 1 条\par 41x Tmt, OURS\par ho\par oe       45\% 70 , Qz ISx0-x) tIo =他-下\par 0 oot            To M tJox J-\%) t15x U-1)" 09\par Me -IS DS + 154-2553\par vat ' 1 y\par oe:\par FEAL AR-G OE 45/359 了                               S75. \textbackslash{}O\par fA 5a £2 FRED\par \par 【第 4 页，来源：ocr】\par ELA     =\par ROT OD ICHT B83\par a by @ A dod Ke A424\par tS hey, 0£\% 41.             (敌人PDM ， las\par seis, 2 NI       (-47A+NX-TYAOronr 24\par te He iG e48 bnd.\par 218\par i\par ad. Cw)             uw (or)\par (c) Sah\par BE going,\par pv     APA TAG F-32 Jott xbs , EzJol\par a         ‘Geo kth =, fez lokr.\par bem 'f          fy\par WACK GRA af ACAGAAO, A Mc: ol,\par (72d me\par Jo\par \#5 gC KHE, Boy: 0 , och "\$44 Ge HA tur finch £8 nly Hd\par Et Bho ty tox JeJobwm\par eeideeting,        bce Ar a1 84 4\par Jo                                 po.\par cokes\par 一全SEE        一二\par 上                                       lo     ,\par mg] Cierny                  Bi uw)\par 141.7/\par \textbackslash{}\textbackslash{} FEL ARS WEBS 46/359 页                            aes\par BAe FRED\par \par 【第 5 页，来源：ocr】\par NLA “乒\par 入币(目1981387888\par io)       Pua\par A pitty                            '\par I,          VA\par WYEKEA dxMaco, orxbxS+9- Glo <0, (95: [uu\par §-0,  ith sarnh ， Fas obbw.\par cha GAG HA, ATONE\par h he gict Rips.\par rs  Lig \#<m,-0, Meg tpx2- 02x27] 20, Hog. 2 oe ghollg po\par P    = 24 bo     \#\par t   a                             ay\par acy esa tt Meg = Mbt lEIEH) 2 tye 3\par hb\par 全  中\par Miac kom)\par OS Be HA HA Co Ge RE ER AA 02 hh,\par ,    x hi                   |\par aa) Ww)\par Me\par 小应土木考研第 47/359 页                            “Sn ip\par iA ae FIRED\par SLRS kee\par \par 【第 6 页，来源：ocr】\par 小让土木考研 语     =\par 去久生上1981387888\par (@)             solv\par 站                 Lom\par An             +\par Ni                 Fe\par WHEGAKEA, HAL pe KE, Wd IMero, b+20+20\%2 - \& X720 下zi\par 3570,  ht -万>o ， ficzly.\par ha Hipee\par Wine 20, wh,       \_ £20, Hy -2 ham\par hg\par bay          on       670,  Qé 4220, Qeg > -2 ky.\par DABS CA KE keh ca FHM mid, anihty 9d (9 0.\par do\par p=\par mig] Chom)            7) chy)\par 中\par I, cpa\par v\par 74          te.\par Whar egg aemaro, E-Ja-4-0.d0 20, Edm\par 352°,  Ar e- Eco,  fiz 了9\par 由 OCP F\par 4     所Mozl Mee tea\par rs\par gig cSt, 07 88 mf 09 08\par +a:                     4a\par ma                 多\par Ne wt '; y\par 小鹿士木考研第 48/359 页                            “Spike\par iA eee FRED\par \par 【第 7 页，来源：ocr】\par “cece, EEOAOATT\par .            IWAN EIOSIO3SZESS\par (a)                 "\par 多\par HH 447GEG, HAL ME, Wsqzo, plas \$x1X2 s)oKeem\par ed    hc facsen,\par WACHAIAY Shr O18 ACR AIH ED Lor 刚 2 asm LOY HY\par wy 08 4             |\par zs\par ACHE H ALE 4 Stor, UR BAL Pa MUME ED 5 rtm\par 5               5\par \{guy                     |\par (b)\par ，  ie A      eee\par f \&.       C\par we betyé 3 0 alhy\% 40 Rice O.\par RH KS A, ray 20, Mg =z Isy Ith <\% kaw,\par £4 Ac He Wh MO, READ Ke BQ\par wy           o ，                  isk, ，，\par | soe          72 i-\par M§) C Kem                                aff Che)\par 小鹿土木考研第 49/359 页                              “Sn Sp\par iA ae FIRED\par Bt-RS ASR\par \par 【第 8 页，来源：ocr】\par I) BLAS         —\par IWAN EIOSIO3SZESS\par (¢)\par A        :       B\par a\par 5\par ie 60M AEG: i - Lurgan La.\par MRI 6, Mo 38 age Le? Cong HD =p\par 8th 6.0 Fe  Bn. tht                       -\par ie OE GAE 7 Herd rng\par Dye = Qsp- Zhe ,  cof AGA 2,8\par £4\par (d)\par Me\par A     0        2 ?zyte .\par 型\par 外\par ip 3 hed TM: 0, *Me- Me \textasciitilde{} fF Azo , §-\% ,\par pes HAY , store, eg = eva\par ZEBS 次M cp z Lal a fy Kt4                .\par tne | 2me .     7\% wih\par we,        ;\par 36,         ag\par 1.v\par OX a nos\par 小鹿土木考研第 50/359 页                        pee\par aN\par Gaba ds' FRED\par Bt-RS ASR\par \par 【第 9 页，来源：ocr】\par 小鹿土木考研       =\par MAC MsOsIOsTCSS\par IO)\par 4          4\par Se\par N              N\par WHERE, dsaaco, tt tx 442420- Ixt7O, foe Il\par Geo! aa cht, fer.\par dB AG C69, Jo br\par \_-  20.\par 4          7      .\par mig Clem\par Qagerbo, Qype ker, PRAY, HAIDA “EER.\par ?L\par 内.\par 及 LE |\par 入       时       fh\par WI ac: éF, mhzo, Aipleg = AC (F594 \#3)\par o ary kA teh.               |\par )\par £49, Qu:\par f      1 By    4        KF\par MY cota att ye = \$f ce Et),\par ; ‘                tHE\par an ar nn nn hae)\par (EEA BB 91/359 页                           “Sr \textasciitilde{}\par 萌喝线上快印\par \par 【第 10 页，来源：ocr】\par 小鹿土木考研                 —\par INA SOsIOSIUSSS\par 中\par r      2h,              9\par °         论\par 有                   hi\par 0 \#1 FAG OX o-0  fax 222\% Ix? 4/ 后 jw\par 了0，  hrf\par Mug      eH), Ree lw\par ‘aii  M0, 4 9 Mey = tlanme   RV ys\par |                                        PO oy deo hh\par BS   ets = a\par ;        图 (ua)\par 了 Race -8 = bby, Ae ARO HH ahem, ad\par :     oe as) cn)\par Vr) el ago ua\par A poetry 多\par G     p      ,\par 1    's\par et seae dg. mg - 72okor\par th CAAWAS ie G\par p\par A pean      ZM: 20， Mo tyr - )oXrx) =2, Meh -— 640 ae\par be 7H Oe tyr -                           ,\par welll wii In              Jo 2 -f 70,  Qck= ops\par ew\par pitti       TH, « Oat chi.\par boa      bag\par wv\par 1.v\par OX LO\par 小鹿土木考研第 52/359 页                            pee\par BA ae ERED\par \par 【第 11 页，来源：ocr】\par 小鹿土木考研                 一\par 站信和微信W93163双858\par hye 9 4\par Wy St.\par hveades WAdahat.\par 190       140\par wok   aa) ce)\par (A)\par WERE, H UL mre, 4/9 -\par c Lg 20 Bg At ep.a2\par )                   “          "A709, Oegz — fn.\par uw b ”\par BOM, lp te\par 有    Moz YB. a 2H.\par Mp\par Ache Higa\par (多\par Ma She            hj 70,   bc 24At Bg =o\par Hed). Hig 4\par Br\par t—>—\_\_                                    4\par =              E               Kf \#1\par zy\par mg.                               多 9A\par vat ' irk\par ON N72,\par ELA WEB 53/359 页                            “Sy “@\par in 5a ERED\par BERS sma\par \par 【第 12 页，来源：ocr】\par ELAS       =\par 未久向全981387888\par ib)        .\par A        Vg       4\par 上 CT rr?’\par hac epee\par A             asm 20, 4\$  2:07\par pt     fi 0 » Bac, tated the,\par 及\par WOES ©) ck BEE\par 1A                    由Imzz0，M .OA 92.\par Que               1 Moe 一4 .\par 954109. 49 oo Gg.\par 一一     gj\par § t= 4.\par mig)                     \%.\par a4,\par N ! irk\par SN Sle\par 小鹿土木考研第 54/359 页                              - Syn S 个\par aa aie FERED\par Bt-RS ASR
\end{enumerate}
\chapter{第4讲：静定梁和刚架}
\begin{tcolorbox}[title=解析总览]
先求支座反力，再分段列剪力和弯矩，最后按突变规律绘制内力图。
\end{tcolorbox}
\section{作业及答案：结构力学基础与技巧班第4次作业及答案（静定梁和刚架）}
\begin{enumerate}[label=\textbf{\arabic*.},leftmargin=2em]
\item \TODO{OCR 自动转写，结构图/公式需对照原 PDF 校对。}\par 【第 1 页，来源：ocr】\par 结力基础班第 4 次作业一一染和刚架\par 粱和刚架作业及答案
\item \TODO{OCR 自动转写，结构图/公式需对照原 PDF 校对。}\par 2-1 图2-60 所示结构截面A 的弯矩〈以下侧受拉为正) BC), CRACKS\par 2008)\par A. M                       B. —M\par C. —2M                     D. 0
\item \TODO{OCR 自动转写，结构图/公式需对照原 PDF 校对。}\par 2-2 图2- 61 所示结构剪力 Fa一         \_。 (BRK 2010)\par M                            Cc  oo a |\par .     )  \textasciitilde{}    B                      |\par C\par A     B\par E        D                                     E       F                  7
\item \TODO{OCR 自动转写，结构图/公式需对照原 PDF 校对。}\par 2-60                            ;      Al 2-61
\item \TODO{OCR 自动转写，结构图/公式需对照原 PDF 校对。}\par 2-3 图2-62所示刚架中AB 杆的弯矩为  ». BAA.  (河北工业大学\par 2010)
\item \TODO{OCR 自动转写，结构图/公式需对照原 PDF 校对。}\par 2-4 图2-63所示刚架，AD 杆截面D SBMn=\_ (AMISEHAIED. CHT\par 大学 2011)\par 10kN-m                                     q\par 4              30kN    |\par \&                       C\par B\par FA            5kN-m                                                   |\par R                      A            \_            B\par ,—4m\_   8m    ,2m 42m)       |   pa\}, 20 \} 22 4 a\}\par 图2-62 2-63
\item \TODO{OCR 自动转写，结构图/公式需对照原 PDF 校对。}\par 2-5 试作图 2-64 所示结构的弯矩图。(北京科技大学 2008)
\item \TODO{OCR 自动转写，结构图/公式需对照原 PDF 校对。}\par 2-6 作图2-65所示多跨梁的弯矩图和前力图，并给出 AB 杆的受力平衡图。 (北京\par 交通大学 2011)\par 本 ‘ie   >     6kNm    | AGN    10kN/m       i\par B                 E    FOG        B       A\par 2m ptm\} 2m\} oo 2a\} mf 2m     \} 5       : 8x2m\par Al 2-64                                         fa 2-65\par 1.v\par AY . “hi Th\par SS A ie i\par O54, es\par Or. inv D\par fA ee FIRED\par Sit-ReE ASK\par \par 【第 2 页，来源：ocr】
\item \TODO{OCR 自动转写，结构图/公式需对照原 PDF 校对。}\par 2-7 求图2-66 所示结构的内力、  作弯矩图(可不写过程)。(东南大学 2008)
\item \TODO{OCR 自动转写，结构图/公式需对照原 PDF 校对。}\par 2-8 绘出图 2-67 所示结构的弯矩图。(武汉大学 2010)\par Fp      Fp            Fp\par 40kN    “on               40kN/m                                                  |\par \}2m\_\{ 2m\} 2m \{2m | \_""4m\_\}       ‘    人2 ee a\par FA 2-66                           ;                                2-67
\item \TODO{OCR 自动转写，结构图/公式需对照原 PDF 校对。}\par 2-9 ER 2-68 Brae nIZe M图。(福州大学 2008)   1           4     .\par z-10 绘制图 2 - 69 所示结构的弯矩图。(武汉科技大学 ye       G 7\par |     ov ATTN...”\par 4 s   \{oY                     I\textasciitilde{}?       E  lm\par D      E       Z\par el ，                      7                       ‘\par wna |             ke            TF 30\par pp                    3m      3m     3m\par 图2-68               图2-69
\item \TODO{OCR 自动转写，结构图/公式需对照原 PDF 校对。}\par 2-11 绘图2-70所示结构的弯矩图。(重庆大学 2012)                                               i
\item \TODO{OCR 自动转写，结构图/公式需对照原 PDF 校对。}\par 42-12 已知图 2- 71 所示结构的弯矩图，试作前力图和轴力图。已知 FP=10kN, M=\par 40kN +m. (哈尔滨工业大学 2010)\par 40\par 20          B         20\par ql                             ‘ "i        =|\par Fp\par c                              ‘\par \{\par \$24" 4+               ;\par 图之-了0                                  Bl2-7\%
\item \TODO{OCR 自动转写，结构图/公式需对照原 PDF 校对。}\par 2-13 作图2-72所示刚架的弯矩图。(中南大学 2009)
\item \TODO{OCR 自动转写，结构图/公式需对照原 PDF 校对。}\par 2-14 绘出图 2- 73 所示结构的弯矩图。(武汉大学 2010)\par Fe=ql          q                        2kN/m        “s               ,\par 3kN\par 1           7                  |   2m   |    2m   |\par 图2-72                图2-73       vel\par O54, es\par Ori VO\par 戎和锡线上快印\par \par 【第 3 页，来源：ocr】
\item \TODO{OCR 自动转写，结构图/公式需对照原 PDF 校对。}\par 2-15 作图 2 - 74 所示结构的内力图。(北京工业大学 2011)
\item \TODO{OCR 自动转写，结构图/公式需对照原 PDF 校对。}\par 2-16 求图2-75 所示结构的内力，并作弯矩图。(可不写过程) (东南大学 2008)\par 30kN-m                               Fp\par (T               mm                             cs\par 30kN         F                         |\par g                                i :                      a\par Z                            时     8\par p24  8m   p—2n +        2+ a :\par : 2-74                                         fA 2-75
\item \TODO{OCR 自动转写，结构图/公式需对照原 PDF 校对。}\par 2-17 作图 2-76 所示结构的弯矩图。(无须写出分析过程) (东南大学 2007)
\item \TODO{OCR 自动转写，结构图/公式需对照原 PDF 校对。}\par 2-18 作图2-77 所示结构弯矩图。(中南大学 2010)      ee\par 20kN/m                                10kN/m           30kN |              :\par \}     C       g       D    Cc    E      F         7\par Z    .                ;                     co\par 4                     .                       \%         \&\par E     A        B            A       =\&B          H\par ¢—Sm\_,-\_Sm\_\_, 6m   pp em ny\par 图2-76                         Kl 2-7   |
\item \TODO{OCR 自动转写，结构图/公式需对照原 PDF 校对。}\par 2-19 已知结构在集中力和集中力偶矩作用下的弯矩图如图2- 78 所示，其中 BC 段\par 为二次抛物线，试确定结构所受的荷载情况，并绘出 BE 段隔离体受力平衡图。 (北京交通\par 大学 2009)
\item \TODO{OCR 自动转写，结构图/公式需对照原 PDF 校对。}\par 2-20 已知图 2-79 所示静定刚架的弯矩图 〈曲线部分为二次抛物线) ，试绘图标出作\par 用于刚架上的荷载，并作出剪力图和轴力图 〈假设各杆无局部平衡的轴向荷载， 无分布力\par 偶) 。(浙江大学 2012)\par 3qa*                      ;                               :\par 2qa@           q2             : D     C     rE\par B    E     qa                          2  \&\par 2ga        8                 5 5      ,\par \&\par A        D                A          B\par MW图(kN'm)\par |一一                   2m     2m   |   ;
\item \TODO{OCR 自动转写，结构图/公式需对照原 PDF 校对。}\par 2-78                              图2-79\par wlan.\par NTA\par SSS A ie i\par 255, =\par COINS)\par BA ee FIRED\par St-RE Abie\par \par 【第 4 页，来源：ocr】
\item \TODO{OCR 自动转写，结构图/公式需对照原 PDF 校对。}\par 2-21 不经过计算，绘制图 2 - 80 所示结构弯矩图。(武汉科技大学 2009)
\item \TODO{OCR 自动转写，结构图/公式需对照原 PDF 校对。}\par 2-22 试求解图 2 - 81 所示刚架，并绘制弯矩图。(东北大学 2005)\par J                        K\par C\par M                                              M \&\par Fp                Fp\par (GF       D                         E       =)\par 县\par 10KN-m\par vt <3,            H-        A                   B           I\par 3m     3m          Sm    3m    有     3m         3m          3m\par ppm fp tp\par 图2-80                                                         Al 2- Si\par o’ ' 1 /\par ONO\par BA ee FIRED\par Bite iz\par \par 【第 5 页，来源：ocr】\par BE EAR StS\par |          TRAC si931 0371666\par .     i,          ws Nowe\par |       4044     网          D> eyo                             |\par i                                                                             |\par TY\par a                 . 和          Mg= Ws     x                          !\par Puy i\par A           多      下 an                                                      4\par C  fh\par tbo                               党         50,         Pe  |\par AAA人           二【bo kum.                 人3\par aw?           F   MAGE 20 yr 2\% Ge 40 oy       sha\par i\}                                 M px                              和\par 加      as\par fe       ri I,        My. Wx 》一PRx4= 3\% lm                      a\par tte                            Nt =9\par SND)\par iA ae FIRED\par \par 【第 6 页，来源：ocr】\par 东久徽售5981887888\par a\par a              Pops ba 信 Aa x3A\par .         A\textbackslash{}L\textasciitilde{}\par | *             4  和  Fas4a+ 和wx ax       ,\par Tite \&\par :     a  从 2 [yin B10)\par 2\par aud ¥\par ;       “      时       4 ppg CES\par ra          WA a).    4\% wr 4e NS a\par |        \} en .    J Mae 4x  J\par |             fy                    人ww     fa\par |                        6 af     Fup ide 7"\par Re                Sv   Typ             — 加。 生\par 5         Kap Se de Stem\par i        (3               lo ¥  -    AAf=  De\par |     is, 一 y aS      wae  vie\par Ww    M2\par |        99\par ]                                Nie\par ET      px Ante BW\par oo       ky    4\% [vx ie os\par Wh rez ?ovxb= th wy\par (ieee                Ape 4 4\par ,        名\par |                                                                                               41 .7\par EPO\par Si:A* ikSieaK\par \par 【第 7 页，来源：ocr】\par |     Aviso 310371888\par |                      we\par |      ye ib\par A  ri           ;   “     如    4?                                 ,\par a            "     4x 4x zz Rlw\par bo         foley  + Gry 4° Siku              |\par ee ex     =\par Epp     mh ae Wek         Mye\par Al be    \textbackslash{}                 ial idl id 的\par ,                       b,\par |            -          bo                                                  .\par oR,                                     bet ahyp\par liad          Fra        ha              .           7\par A         多          5            fda”    Aa ten                               \}\par f J                            |                                          |\par ft           Maz 2Fpxa- FpxAz Fh,                                      , |\par 2, P                                M =o\par FS                 Ax Ye Dt Ixy                   J     a\par 人\par ff  VS:       Fake hw              :\par [上 ena alee                 :\par pn      it 4  + zy 2二|\par D leo          wv Mts bxI- 244 = 人-= kvm,            7\par .                                                                                          1.v\par NO\par A ae FRED\par \par 【第 8 页，来源：ocr】\par CACM EI 931 O31,80.9\par 和 is,          tof]              ibs -                            |\par wy |v oP  /”     hp   mo stot af\par |        ia      Bp culty              InX aud + Fore zh\par |                  jw     4 shail\par le!   NF     a 本本\par TN\par >a ay  | Mw\par aN      PY Poy vi  B\par |     LA     ，        \_、\par | |     /    fav       Mp= W  , “\par |      A <           各xz- Uns *             |\par |  ‘           i ne  bs. Mg Blt uy =冰\par | i        4o            cs  Vv\par |      vs       0\par feel 22\par 人                         FF二\par |     mo                     中     al |\par joo 中                bw\par 人网\par NO\par 萌哆线上快印\par \par 【第 9 页，来源：ocr】\par EY SBAR Sis\par |        TRAM si9 31031666\par by                                          ,\par C\par A             二人gm\par way”                 “Tay    hye Fay oe                                   !\par a                                                                      |\par 4         加                        国| 本 本\par *\par ma                 b                                                          ,\par 2\par 244,             4\par 2 1中\par Ww      Ap + \textasciitilde{}24|=8\par Nuz 3x4 x2 \%2- he\par b,                     t\par 3\par |                 》\par 7S. \textbackslash{}     -                                 Mazes\par 和          lox 4\par fox k    |                    A eli            4olw\par lm     © “hl                    vo\par ,        U                          Bt.\par 1 .v\par ONAL\par ENO\par ieee FRED\par \par 【第 10 页，来源：ocr】\par RAR Stil\par ARO RU OBIGETOSS\par 40x 4- WK Y= Inx (vy Y=                           |\par vite 一 M fue 4x 4- were \$0 my,\par |    a)        sob    “\par °                                 \}\par |       PN         可\par |     yo\par | |      ) i  (i           f\par |      \_   AL)\par ee:\par . |    »  x 30 ”8.75   ok                 Te |             |\par |  red |            40  8.75 30     |\par ele                              oe               8.75 leat 10\par we\par Ci @\par iA ae FIRED\par Siew nse\par \par 【第 11 页，来源：ocr】\par EEL SBAR STS\par IRM OD IGSTOSS\par 24h, \& 伸         rr          :\par To                      b  be  Bb\textbackslash{}\textasciitilde{}\par fr >\} im     In      fp :      了po [un\par = /Mpa\par a  b   4 |       h           She ag\par 、 Mave Fre Fra.\par AOS Foye H\par “i i\par ‘                                              一>\par al           SA            |\par |                                 L\par 中                              a\par \textbackslash{}4\par joy bz |? tt\par A      a\par 的人\par A ae FRED\par \par 【第 12 页，来源：ocr】\par 未久微信19812587888\par 三Ap        fo x TS 4 Wx Ib                              |\par \_              \& Mo»                                          E\par fe  4       G         4       F        box b+ Wx” TS 4 Wy DYb\par “3                         fyy (=\par ce Fy = ib8 TS lov\par 一个            F   “TfR Feent fe                                         :\par a                Fry     Frxx \{+ Wr by\} 2  ib¥TS\% 6\par |                               5       ce Efe P25 kw\par BLS      》  \textbackslash{}  Fox + bot \{vx 4\par ‘   a hye 1S\par a              a Poxz 站 kw\par 10S                       Neve boy 3t IFRS -      -\par fxn ir" 3bv                         加\par M (him) “                                             “8\par to Pee ry fee by  UT He                    7\par oF           AE FL\par ,                           GT fag I 4¥ 4oy. 4= 4p,          ，             \_\par is          Yasin «Een a5 tay Vs  4               :\par \_ BAA\par °              Foyx oe low 和2                        ed\par 一人人          a le            |              I\par 个    “a    Foxx 4+ Joy Y= Fo + 2 \textasciitilde{} WUy2             Pe\par fe\par 5        SS Wye Dw                                          :\par >" (oy 4's Ww.                           ;\par ae 1 。/\par OPO\par fA eZ FIRED\par \par 【第 13 页，来源：ocr】\par 后绢六菠公众呈LA\par TRAMs Os IO3TSCS\par 2414,     WY\par ave          < 4A               oy\par Ly                4X fu w= + waz A                  a\par |     a               wfe bY\par |                     =>     ¥      r\par |              Er                  “      ii a                         \textasciitilde{} ap")\par |                              wae\par \}                    \textbackslash{}                os  Wy\par 7M,\par \$e fe 42 XS 下 fe b kw\par |     和\par Nye Yew\par |                           Pour 4-24 2\par Ae        b     —> TR = ASE                                                  |\par f | 4\par i | Ra图GD)  机  "9 t ‘  FyR(KN)  ae a\par ;                                                                                               ae ee)\par CRO)\par 萌喝线上快印\par \par 【第 14 页，来源：ocr】\par RAR Stil\par RENE ODIGBISSS
\item \TODO{OCR 自动转写，结构图/公式需对照原 PDF 校对。}\par 2-2" 、             ib          lo              /| \textbackslash{}   Fee  By ut ¥y KIA H r\%\par a ee avy aitep, meaty uy\par 2722,\par mr       MM                                  AN                                              |\par sv                     *\par a ae rail oof                             ,\par AF Yo Hap top\par oft                 f   DY 7H OS RTH,                       .\par For ADU Rep\par REfp         B    Lp  W\textbackslash{}t= fox 六fr 2 af.\par aff               3h                                           |\par “Sp\par 萌喝线上快印
\end{enumerate}
\chapter{第5讲：静定桁架}
\begin{tcolorbox}[title=解析总览]
优先判别零杆，再用节点法或截面法求目标杆件轴力。
\end{tcolorbox}
\section{作业及答案：结构力学基础与技巧班第5次作业及答案（静定桁架）}
\begin{enumerate}[label=\textbf{\arabic*.},leftmargin=2em]
\item \TODO{OCR 自动转写，结构图/公式需对照原 PDF 校对。}\par 【第 1 页，来源：ocr】\par a REAPS Ra ES BRAS eB an OT ro\par 9   ，\par BS —   S\par Fp       .\par a   |  a ¥  a                                                 .\par (a)\par Fl 2-139\par 六、练习题
\item \TODO{OCR 自动转写，结构图/公式需对照原 PDF 校对。}\par 2-23 图2-140 所示柏架中零杆的数目为〈不包括支座链杆) CD. (中南大学\par 2010)\par A. 6根                      B. 7根\par C. 8根                      D. 9根
\item \TODO{OCR 自动转写，结构图/公式需对照原 PDF 校对。}\par 2-24 计算图2- 141 所示梅架中1、2 杆件的轴力。(重庆大学 2010)\par I   I\par 图 2-140               fl 2-141                 ney\par |                                                       ON,
\item \TODO{OCR 自动转写，结构图/公式需对照原 PDF 校对。}\par 2-25 RE 2 - 142 所示结构中指定杆件 1, 2 的轴力。(广州大学 2012)       NE
\item \TODO{OCR 自动转写，结构图/公式需对照原 PDF 校对。}\par 2-26 RE 2 - 143 所示结构中指定杆件 1、2 的轴力。(广州大学 2008)      2 Ee:\par O70\par 78           A ee FRED\par BERS Asie\par \par 【第 2 页，来源：ocr】\par B\par :ax 。              ;        \textasciitilde{}\par I L- |    4       D   F     《\par dy           by             Fp        Fp\par 了   I   i            7    i    i\par rs rs        人 人 rs |        oO\par Al 2-142                Kl 2-143         |
\item \TODO{OCR 自动转写，结构图/公式需对照原 PDF 校对。}\par 2-27 求图 2 - 144 PRAHA 1, 2 杆的轴力。(西南交通大学 2008)
\item \TODO{OCR 自动转写，结构图/公式需对照原 PDF 校对。}\par 2-28 作图 2 - 145 所示枯架中 a, b,c 杆轴力。(西南交通大学 2010)\par Fp\par |p        I>          Nes\par . ,—4 a |  papa a\par 图2- 144                  图2-145      .
\item \TODO{OCR 自动转写，结构图/公式需对照原 PDF 校对。}\par 2-29 BORA 2 - 146 所示柏架中 w、8 杆件的内力。(东北大学 2007)
\item \TODO{OCR 自动转写，结构图/公式需对照原 PDF 校对。}\par 2-30 求图 2 - 147 所示桥架中 w、杆件的内力。(北京航空航天大学 2008)\par : Fp=1  0          |\par TASHA\par 2d   2d        6x1         |\par 图2-146                     图2-147                       1\par Ne ‘ ;e 4
\item \TODO{OCR 自动转写，结构图/公式需对照原 PDF 校对。}\par 2-31 计算图 2- 148 ARH a, OHA. (BRAS 2008)           RONG
\item \TODO{OCR 自动转写，结构图/公式需对照原 PDF 校对。}\par 2-32 求图 2-149 所示顶架中 1, 2 杆件的轴力。(西南交通大学 2007)         Tit oe\par 2 es\par |            o/iw@\par 7戎锡线上快印\par \par 【第 3 页，来源：ocr】\par af FREAS SRN ES BRA Ae DLS FATWA\par 了  °  F                   t L交和、人|\par 1 |一: 证  *  a  f  a  ;\par Al 2-148                    fl 2-149\par 七、练习题答案
\item \TODO{OCR 自动转写，结构图/公式需对照原 PDF 校对。}\par 2-23 C,
\item \TODO{OCR 自动转写，结构图/公式需对照原 PDF 校对。}\par 2-24 Fy =V2Fp/2; Fy =—Fp/2.
\item \TODO{OCR 自动转写，结构图/公式需对照原 PDF 校对。}\par 2-25 Fw=—50kN; Fy=—50/2/3kN,
\item \TODO{OCR 自动转写，结构图/公式需对照原 PDF 校对。}\par 2-27 Fu= —J2Fp; Fe=0.
\item \TODO{OCR 自动转写，结构图/公式需对照原 PDF 校对。}\par 2-28 Fu = —(2Fp/2; Fy =V2Fp/2; Fy.=Fp/2, fem: 先求支座反力，再依次取\par 截面IT-工、I-I、亚-焉，见图2-150。由工-工截面求 Fw OAPs 轴列投影方程)，\par ” 由开-开截面求 Fe，由正-亚截面求 Fu。
\item \TODO{OCR 自动转写，结构图/公式需对照原 PDF 校对。}\par 2-29 Fu一一Fp; Fe一V2Fp/2。
\item \TODO{OCR 自动转写，结构图/公式需对照原 PDF 校对。}\par 2-30 Fu=0; Fe=0.5。提示: 集中力右侧和悬臂部分的杆均为零杆，去掉零杆后\par 结构变为图 2 - 151，由对称性结论可以判断出 ao 杆为零杆，其余零杆也示于图中，20 杆的轴\par 力由结点法求出。\par Felony\par el\par i    NS                |fe\par ，      I-       f<-1                0\par ce    ee\par Fp            I               0\par 图2- 150                 图2-151               、、11y\par ON,\par .<     “7 o
\item \TODO{OCR 自动转写，结构图/公式需对照原 PDF 校对。}\par 2-32 Fu —/8Fr/2s Fe —Feo SR: 先判断出 A 点连接的两根斜杆为零桩:2 CO STS\par 再用结点法或截面法。                                            OPENS)\par . 1 -\par \textasciitilde{} 80                                                      萌锡线上快印
\end{enumerate}
\chapter{第6讲：组合结构}
\begin{tcolorbox}[title=解析总览]
先拆分组合结构的梁、桁架和刚架部分，再保证连接处作用力平衡。
\end{tcolorbox}
\section{作业及答案：结构力学基础与技巧班第6次作业及答案（组合结构）}
\begin{enumerate}[label=\textbf{\arabic*.},leftmargin=2em]
\item \TODO{OCR 自动转写，结构图/公式需对照原 PDF 校对。}\par 【第 1 页，来源：ocr】\par aR ee Ree Os Shae eR IT Pa adas\par 根据结构与荷载特点，尽量利用简捷的方法求出结构的弯矩图。(武汉理工大学 2010)\par Fp     Fp\par Y A   B   F   C\_\par i if 4 12 a M I  .\par (a)\par 图2-162\par 三、练习题
\item \TODO{OCR 自动转写，结构图/公式需对照原 PDF 校对。}\par 2-33 图2-163 所示结构中杆AB 的轴力为(  )。(了哈尔滨工业大学 2011)\par A. 0                        B. 一0.894 Fp\par C. —0. 5Fp                  D. —Fp
\item \TODO{OCR 自动转写，结构图/公式需对照原 PDF 校对。}\par 2-34 对图 2-164 所示结构，作受弯杆件的弯矩图，算出各链杆的轴力 〈标注在杆件\par i). (北京交通大学 2010)\par 9\par 、       a    a |B a   ‘a\par Al 2-163                                Kl 2-164
\item \TODO{OCR 自动转写，结构图/公式需对照原 PDF 校对。}\par 2-35 对图2- 165 所示结构作弯和矩图，并求杆件 1 的轴力。(广州大学 2011)
\item \TODO{OCR 自动转写，结构图/公式需对照原 PDF 校对。}\par 2-36 试作图 2-166 所示结构的弯矩图、剪力图、轴力图。(哈尔滨工业大学 2010) 、 Ah
\item \TODO{OCR 自动转写，结构图/公式需对照原 PDF 校对。}\par 2-37 作图2- 167 所示结构的灾矩图。(重庆大学 2010)                   LONE
\item \TODO{OCR 自动转写，结构图/公式需对照原 PDF 校对。}\par 2-38 fe 2-168 RAMA, RAY 2012)               Trt fz.
\item \TODO{OCR 自动转写，结构图/公式需对照原 PDF 校对。}\par 2-39 ER 2-169 Pra, FRAT EHR. CINK 2007)  ey    eas\par .                                    LON   .\par O70\par 86                                                          BA eee FIRED\par St-RS ksi\par \par 【第 2 页，来源：ocr】\par a EY Soro tt is\par Fp
\item \TODO{OCR 自动转写，结构图/公式需对照原 PDF 校对。}\par 0-0          ;           :           G\par q\par 20kN/m    20KN | \&                   “      \textasciitilde{}\par PTL Tit                       /\par 1            \textasciitilde{}                       7                    四\par 一          \&                              \textasciitilde{}\par 1     1    \{     4m    f 2m |         |    1   |    7   l   |\par 图 2- 165                 图2-166           图2-167\par 本    Fp\par 2kN/m                Fo    Fel      ;\par A                 Cc        T       VA      .\par B                            \textasciitilde{}\par §\par F                        wa          -\par §                                        \textasciitilde{}\par ;                               D     ;\par 4m  :   4m       本     |一全\par 图2-168          图2-169\par 四、练习题答案                  -
\item \TODO{OCR 自动转写，结构图/公式需对照原 PDF 校对。}\par 2-33 C,                                                    |
\item \TODO{OCR 自动转写，结构图/公式需对照原 PDF 校对。}\par 2-34 答案见图 2-170,
\item \TODO{OCR 自动转写，结构图/公式需对照原 PDF 校对。}\par 2-35  Fu = —2ql; 2s 6 PR] OF] 2-171.
\item \TODO{OCR 自动转写，结构图/公式需对照原 PDF 校对。}\par 2-36 弯和矩图、剪力图、轴力图如图2- 172 (a. (b). (c) 所示。\par a\par a    a\par ©   0   人\par \_ ga        \_ga             .       P\par 六       Q              ——  0          :\par . 2g?     1\par |                           :              Cae 1 ' iY\par O—O                    aw,                                          (0)  . Ae  (0)\par O                           O    =   人 ;            NS  “Se\par ME Fv                        MB                   >:         oa\par | 2-170 2-171                         es         it\par oF).  ws. 、\par o, bare .\par O75 VO\par 878 ES FIRED\par \par 【第 3 页，来源：ocr】\par of PRESS RNS OS BRE AR | 2 LGR [(FanAadaar\par 140 |       ,                           :        280\par \textbackslash{}                   3                                  =\par pa enna       |       一全-\par 50\par iw           300\par 2      4图(kNm)              2      Fo图(kN)     及图kN)\par (a)                           (b)                       (c)\par BI 2-172
\item \TODO{OCR 自动转写，结构图/公式需对照原 PDF 校对。}\par 2-37 SRA 2-173 所示。
\item \TODO{OCR 自动转写，结构图/公式需对照原 PDF 校对。}\par 2-38 弯矩图如图2-174 所示。
\item \TODO{OCR 自动转写，结构图/公式需对照原 PDF 校对。}\par 2-39 弯矩图如图 2- 175 Bra.\par G                                                   D\par F           4g       :       9                p78\par Ver\par Fpl     E                  \textbackslash{}      32\par oy \textbackslash{}24                            :4\par ue 学          na         ia Sey\par 图2- 173                 图2-174              图2- 175\par 第四节 三 Be Ht\par 一、拱的受力特点\par (1) 在竖向荷载作用下产生水平推力。\par (2) 因为水平推力的存在，使三铵拱的弯矩比相应简支梁的弯矩小。\par |        (3) 在竖向荷载作用下，梁的截面没有轴力，而拱的轴力较大，且一般为压力，因此，\par 拱比梁更便于利用抗压性能好而抗拉性能差的材料。                                        、、 1 ' i! y\par No\} tes O,\par .       =, =Reeete Ra PTR AXk                         shes\par 1. 支座反力公式                                                          oo      oe\par 7 soe tN   .\par (1) FH CHL 2-176 〈a)]                                                    CONN)\par 88                                                                          BA eee FIRED\par St-RS isi
\end{enumerate}
\chapter{第7讲：三铰拱}
\begin{tcolorbox}[title=解析总览]
三铰拱题利用中铰弯矩为零求水平推力，再求截面内力。
\end{tcolorbox}
\section{作业及答案：结构力学基础与技巧班第7次作业及答案（三铰拱）}
\begin{enumerate}[label=\textbf{\arabic*.},leftmargin=2em]
\item \TODO{OCR 自动转写，结构图/公式需对照原 PDF 校对。}\par 【第 1 页，来源：ocr】\par Be RAPES EP tty\par 五、练习题
\item \TODO{OCR 自动转写，结构图/公式需对照原 PDF 校对。}\par 2-40 图2- 188 所示三铵拱截面 K BAH       (下拉为正)，拱轴线方程为他\par (lz—z’), (广西大学 2010)
\item \TODO{OCR 自动转写，结构图/公式需对照原 PDF 校对。}\par 2-41 Al 2-189 PRESEN KERERABBL? 〈武汉大学 2010)\par q=20kN/m\par |                 20RN    20kN\par 10m          |\par 6m            6m      \{   3m   3m       6m       \}\par A 2-188 2-189
\item \TODO{OCR 自动转写，结构图/公式需对照原 PDF 校对。}\par 2-42 图 2 - 190 所示结构中，链杆 AB 的轴力为          。(以受拉为正) (中南大学\par 2010)
\item \TODO{OCR 自动转写，结构图/公式需对照原 PDF 校对。}\par 2-43 图 2- 191 PRES, BERATED y= FF 2-2), JOP f=4m, 1=16m,\par 则截面忆弯和矩 Mp等于 x RA DWAFoOSF\_\_\_. CHEK 2012)\par A                    B       5                         °           D           \&\par 2                                            g é\par pla, 2                 上-和 dm tm\par Al 2-190                                     |         2-191\par 六、练习题答案
\item \TODO{OCR 自动转写，结构图/公式需对照原 PDF 校对。}\par 2-40 13.33kKN . m。
\item \TODO{OCR 自动转写，结构图/公式需对照原 PDF 校对。}\par 2-41 45kN,
\item \TODO{OCR 自动转写，结构图/公式需对照原 PDF 校对。}\par 2-42 —Fp UB).
\item \TODO{OCR 自动转写，结构图/公式需对照原 PDF 校对。}\par 2-43 4q (上侧受拉); 0.                                        、\par Ve:\par 253, oe\par |                                                   ofA S242 FRED\par Bi-ReE Asia
\end{enumerate}
\chapter{第8讲：静定结构影响线}
\begin{tcolorbox}[title=解析总览]
影响线题可用静力法或机动法，关键是单位移动荷载位置与正负号。
\end{tcolorbox}
\section{作业及答案：结构力学基础与技巧班第8次作业及答案（静定结构影响线）}
\begin{enumerate}[label=\textbf{\arabic*.},leftmargin=2em]
\item \TODO{OCR 自动转写，结构图/公式需对照原 PDF 校对。}\par 【第 1 页，来源：pdf-text】\par 第三章\par 静定结构的影响线\par 扣\par ＋\par N\par =\par F\par 沪max=\par lOXO. 5 =25kN。\par 五、练习题\par 图3-48所示结构在任意长度可移动荷载20kN/m作用下，截面K的剪力最大\par 值FQKmax 为\par 。(天津大学2010)
\item \TODO{OCR 自动转写，结构图/公式需对照原 PDF 校对。}\par 3-2求图3-48 所示结构在任意长度可动均布荷载20kN/m作用下截面K的最大弯\par 矩MKmax l。(华南理工大学2013)
\item \TODO{OCR 自动转写，结构图/公式需对照原 PDF 校对。}\par 3-3\par 图3-49所示多跨梁，截面K的剪力影响线在K点的竖标为YK左\par ＝\par y店; =\par ; 在图示移动荷载作用下，截面K的剪力最大值为\par 。(浙江大学2009)
\item \TODO{OCR 自动转写，结构图/公式需对照原 PDF 校对。}\par 3-1\par 123\par \par 【第 2 页，来源：ocr】\par al 研究生入学考试辅导丛书 结构力学〈第二版)          TS oO LNIESLLSE \$7 /CeXeXe)\par .     -                    a | jor\par 20KN/m\par |  4m  i   8m 2m |   6m   |     |  5m  |   8m  |, 2m | 2m ) 8m\par 图3-48                         '     图3-49
\item \TODO{OCR 自动转写，结构图/公式需对照原 PDF 校对。}\par 3-4 用机动法作图 3-50 所示多跨静定梁中 Ma 、Fas 、ME 、M 的影响线。 (湖南大\par 学 2008)    |
\item \TODO{OCR 自动转写，结构图/公式需对照原 PDF 校对。}\par 3-5 作图 3 - 51 所示结构的 Foxe. Fox. Mx 〈下侧受拉为正) 影响线。 (广州大学\par 2011)\par Fpl    .      .                     Fpl\par 3m   [1m | tm | 1m | Im  2m  | 1m|  2m    2m  2m  2m
\item \TODO{OCR 自动转写，结构图/公式需对照原 PDF 校对。}\par 3-50      .                   图3-51
\item \TODO{OCR 自动转写，结构图/公式需对照原 PDF 校对。}\par 3-6  用静力法作图 3-52 所示结构的 Me、Fas 、MAa影响线。Mc以杆 DC 下侧受拉为\par TE; Ma以AK 杆端下侧受拉为正。(Fe一1在 FDCH 段上移动) (青岛理工大学 2012)
\item \TODO{OCR 自动转写，结构图/公式需对照原 PDF 校对。}\par 3-7 ER 3-53 所示结构的 ME、MFr影响线。(华南理工大学 2008)\par Fp=1      :\par .                       A               E\par x [Fl                                                      呈\par F一7    GC        H    ,                                 只\par | 2 | 24 4 22 -| 2 4    p—2m | 2m, 2m, 2m 4, \_2m\_\par .       图3-52                 图3-53
\item \TODO{OCR 自动转写，结构图/公式需对照原 PDF 校对。}\par 3-8 移动荷载 Pe—1 PRB, MEM 3 - 54 PTR EAD Fy 和 Py\par 线。(河北工业大学 2008)\par NEES Kf\par .       ‘          2mx6-12m     i  .\par 图3-54\par aes
\item \TODO{OCR 自动转写，结构图/公式需对照原 PDF 校对。}\par 3-9 作图3-55 SRARZA Fra. Fu. Fre. Fes Renee. Cu Arte EF RAGES NO,
\item \TODO{OCR 自动转写，结构图/公式需对照原 PDF 校对。}\par 3-10 (ERI 3-56 RARE Fea. Feo. Mob. INK 2012) ZN\par O70\par 24                         |                              iA ee FRED\par \par 【第 3 页，来源：ocr】\par SS) CET Se Voy\par x       Fp=1                 .      、\par NA\par |            6xa\par Al 3-55
\item \TODO{OCR 自动转写，结构图/公式需对照原 PDF 校对。}\par 3-11 图3-57所示结构，单位荷载 PP=1L AAA, HB. C. Diam. (1) 作轴 -\par 力 Fog 、剪力 Faee和Mr的影响线。 (2) 若 CD 杆受到图示荷载作用，利用影响线计算\par Fype Fl Mac. GRINFE LAS 2011) ©\par ‘  Fe=l ”“      ;      > 4KN/m\par t=\}\par *                          g\par A           B     、\par C                    A               E\par 4m  |  4m  |  4m  |  4m |  3m  ,  3m  ,  3m  "    .\par 图3- 56       图3-57
\item \TODO{OCR 自动转写，结构图/公式需对照原 PDF 校对。}\par 3-12 BOR 3-58 所示结构内力影响线 Ma Fm, For (单位荷载 已一1在 DF段移\par 动) 。(青岛理工大学 2008)
\item \TODO{OCR 自动转写，结构图/公式需对照原 PDF 校对。}\par 3-13 图 3- 59 所示结构， Fp=196 BCD 移动，试作出截面 K SH 〈下侧受拉为正)\par 和轴力的影响线。(浙江大学2012)\par Fp=1               |\par D     E.        F         B IF »=1\par B          Cc                                   一                  Lo\par A                      4  K  mC OD\par |  了7  |  7  |  7  |  1         | 1/2 | 1/2 |  I   \}\par 图3-58                      图3-59
\item \TODO{OCR 自动转写，结构图/公式需对照原 PDF 校对。}\par 3-14 用静力法作图 3 - 60 所示组合结构的 Ms 、Fn 、Fie的影响线。(青岛理工大学\par 2009)
\item \TODO{OCR 自动转写，结构图/公式需对照原 PDF 校对。}\par 3-15 图3-61 i mA—-T=EBH: (1) 设方向向下单位集中力 Fp=1 可以在拱上沿x\par 轴方向移动，求作支座 A 水平反力 Fa的影响线 (z轴以A 为原点，以向右为正向) 。 (2)\par 设拱上布满方向向下、沿 z 轴方向均匀分布的荷载g，试利用 CL) 中所求影响线， 求解水，，, ，\par FRA Fn的值。(武汉理工大学 2008)             SN
\item \TODO{OCR 自动转写，结构图/公式需对照原 PDF 校对。}\par 3-16 (1) 作图 3 - 62 RA DH Fr MK 点弯矩 Mr的影响线; arse GG zo:\par 示移动荷载组作用下 Mk的最大值。(要考虑荷载掉头) (西南交通大学 2012)     人\par seis’ Swen\par BER iS\par \par 【第 4 页，来源：ocr】\par a 研究生入学考试辅导丛书 结构力学〈第二版)           12 Y< SAG AIRES Eee\par x   Fp=l\par |    |                 Fe=1\par E        D                         x\par C                                          4   C\par \&                   y                         \textasciitilde{}\par A                     B              \_\_\_, 4  \_              D\par | -一阵 | im                人   |天一|一\par 图3-60                            .               图3-61\par oo 上 poe\par Fol                       im 2m |\par A        K     B      C    D    E     F\par p—6m\_, 4m\} 4m 4 4m 4 4m 4 4m\_y
\item \TODO{OCR 自动转写，结构图/公式需对照原 PDF 校对。}\par 3-62                               :\par 1./\par Ne   ry\par iT\par CA NS)\par 12                                                                                    Zi\par 6                                                                             萌锡线上快印\par \par 【第 5 页，来源：ocr】\par FERRARI\par BARS         REE ODIGHUSES\par 4H.          |                                               |\par |                        ym                                        (\par Bhs (Hele = pk\par |      关公\par |      uh\par |                                   2m     fm\par |       | kms > xx Wee Yo bem                                       .\par 44,\par |                       Benne Sxl += Tol                          ,\par FAR   \%  \%                            |                   |\par |      oN\par 的\par ci                            荫喝线止快印\par \par 【第 6 页，来源：ocr】\par BE RBAR Sis\par |               |                              TRACE g3 1037888\par :      OD mM                        a\par @ My              1\par oO Pete\par |                                                                               |\par ® Me                                                  \_\par 3-b.               PA\par OMe      A          aT?                               7\par vey O24                                                                                \&,\par M br =  at w= RAM) 2 L (ta-ey—FAt Xe 2A SKA\par Teg              Ode:\par 一    Va          ieee FRED |\par \par 【第 7 页，来源：ocr】\par HE SBAR Sis\par kak\par Pye ao                                                    !\par |    xzmco , Foyy ude Ix[x-\$4)=0, y= co 四\par —«® fee             Be         ——  |\par |    ox esa fv       Fh            Es m =\par pee                                    m\par Ta Pee a                         Fd.        |\par @) of Xu a,   Rh= ax  mez  这 “RA -  1(44\textasciitilde{}x) =a                 |\par : |                                        4a\par we 和\par :                                                        GENO\}\par 二                             荫喝线止快印\par \par 【第 8 页，来源：ocr】\par —      PRRBARS
\item \TODO{OCR 自动转写，结构图/公式需对照原 PDF 校对。}\par 9-8,  vi 1ia\textbackslash{}\}      A  Jd        IRAN OB TBE B88\par fe   :   人   人     人   0      n”   TR\par PTS BE\par |                            ta N i =              |\par |                        5        =\par Om Daal  en 2 i\par eal VL 7力克] ws]  < +)     i     ye            ;\par 34                                            |\par Cay                                    |      7\par VK ph A, Via SN\par fo ”             ate           本\par je\par -AN\par ph Uk h TPS)     2\par 二                      tw2efe, pH        。\par 4 EDS. esY Cele\par Ae\par fA See ERED\par \par 【第 9 页，来源：ocr】\par a          PRRAAR SN\par KZ 东久向全1981387888\par tall\par CL /\par WN     0     ZN/\par Fa          b              °\par 加                          Foo] HEE MI RSY,\par ha: Fe , Fs FS \%) 04,\par 号            HAR), Fate Fs\par Ae ie\par “Sy oO\par BA 22 CRED\par \par 【第 10 页，来源：pdf-text】\par 当单位力作用在下弦时\par (1)\par RA\par F 杆轴力影响线\par 与作用在上弦相同，可直接画出单位力作用在下弦时\par RA\par F 影响线，如图(a)所示。\par (2)1 杆轴力影响线\par 采用联合法，断开1 杆[见图(b)]。令断口两端发生单位相对位移，AC、DB 段影响线应分别保\par 持为直线，取A、C、D 和B 为关键点。当单位力作用在A 点和B 点时，力通过竖杆直接传递给大地了，\par 所以yA=yB=0。当单位力作用在C 点时，如图(c)，求出支座反力，取Ⅰ-Ⅰ截面以右部分隔离体，根\par 据竖向力平衡可求出\par 1\par 2\par 6\par N =\par ；当断开1 杆后，链杆IJ 和CD 相当于定向节点，因此AC 和BD 两段的\par 影响线应保持平行，根据AC 段影响线，可以求出单位力作用在D 点时\par 1\par 2 2\par 3\par N =\par 。连接所有关键点竖\par 标，可得1 杆轴力影响线，如图(d)所示。\par (3)2 杆轴力影响线(9 分)\par 当单位力作用在下弦时，2 杆始终为0 杆，影响线而(e )所示\par \par 【第 11 页，来源：pdf-text】\par (4)3 杆轴力影响线\par 采用联合法，断开3 杆[见图(f)]。令断口两端发生单位相对位移，AE、BF 段影响线应分别保\par 持为直线，取A、E、B 和F 为关键点。当单位力作用在A 点和B 点时，力通过竖杆直接传递给大地了，\par 所以yA=yB=0。当单位力作用在E 点时，如图(g)，求出支座反力，取Ⅰ-Ⅰ截面以右部分隔离体，根\par 据竖向力平衡可求出\par 3\par 2\par 2\par N =\par ；当断开3 杆后，链杆KL 和EF 相当于定向节点，因此AE 和BF 两段的\par 影响线应保持平行，根据AE 段影响线，可以求出单位力作用在F 点时\par 3\par 2\par 3\par N =\par 。连接所有关键点竖\par 标，可得3 杆轴力影响线，如图(h)所示。\par \par 【第 12 页，来源：ocr】\par BE BARS GS]\par ,     to,\par OP 一一           ihn |, BOFTT AH         |\par 7                 =\par aa         g        >\par .                        |        中                                   、、41.y\par ae Z\par SRO\par |                                     fee FIRED\par \par 【第 13 页，来源：ocr】\par |     —-—\_-—\_SSRRAaReind —\par a       术欠向信2991587888\par b             b             ,       Br Bek 4\par DEA YE, HAS fs RB        i\par Po \& Mazo                            7\par yp\par Mae 2 Pey                     7\par ¥                                                          :\par AO\par aA eee ERED\par \par 【第 14 页，来源：pdf-text】
\item \TODO{OCR 自动转写，结构图/公式需对照原 PDF 校对。}\par 3-11、(1)\par 解：①\par NDE\par F\par 影响线\par 采用联合法，取B、C、D 为关键点，其中当单位力作用在B、C 点时DE 杆轴力为0.当单位力作用在D\par 点时如图(a)所示，根据结点D 竖向力平衡可得\par 5\par 4\par NDE\par F\par = −\par 。综合以上计算，画出影响线如图(b)所示。\par ②\par QBC\par F\par 影响线\par 用机动法。去掉相应的约束，画出虚位移图[见图(c)]，令\par QBC\par F\par 对应的位移等于1，求出影响线的竖标\par 后可得到BC 杆B 端截面的剪力影响线如图(d)所示。\par ③\par BC\par M\par 影响线\par 用机动法。去掉相应的约束，画出虚位移图[见图(e)]，令\par BC\par M\par 对应的位移等于1，求出影响线的竖标\par 后可得到BC 杆B 端截面的弯矩影响线如图(f)所示。\par (2)图示荷载作用下，可分别求得DE 杆轴力和\par BC\par M\par 为\par 1\par 5\par 15\par 3\par 4\par 2\par 4\par 2\par NDE\par F\par KN\par = −\par \textbackslash{}\textbackslash{}times \textbackslash{}\textbackslash{}times\par \textbackslash{}\textbackslash{}times\par = −\par 1 3 3 4\par 18\par 2\par BC\par M\par KN m\par = −\par \textbackslash{}\textbackslash{}times \textbackslash{}\textbackslash{}times \textbackslash{}\textbackslash{}times\par = −\par \textbackslash{}\textbackslash{}cdot\par \par 【第 15 页，来源：ocr】\par PPAR\par Se:    |\par ,                                   .                                                         GENO\par 萌喝线上快印\par \par 【第 16 页，来源：ocr】\par -    SASRARSION\par ce tht            未久向信1951587888\par 39. Feb dik, Mes\par 1 A> 8 \textasciitilde{}     B        At   yey\par ae\par k    c      D\par se\par as i ca         |\par ha 1ERA    4: Tis ee\par b        cl       p’         有\par \textasciitilde{}\textasciitilde{}                 全\par 是            iL       了\par mening >           Foot. By oft) 4X                   -\par aie, \textbackslash{}f 4         4CXL3m\par 1          4        HII 4, mo\par ,                             M20            ;\par |                    f         Hy B.     Netra 2 Ka gf= 5               (       |\par Me 53      = |        |\par 和CC Agz=X                           mo Ree 4x        :\par BA  cA  an        一\par Q 人\par rad                 \_ wets 2   r       部\par A\textbackslash{}cvz                     -\par |                                  ,   Ay Hoh sx                 3\par |                                                      Noe\par 萌喝线上快印\par \par 【第 17 页，来源：ocr】\par 全刘\par AAS                          |\par \{)      jer | C              aint t\par |                 AN 30          ，\par |           |                   N\textbackslash{}=0                  ，\par ¥\par aff      \textbackslash{}          4\par .      wey       y+ fox te Fy-y         ;\par 5   AM +  Font =   Ka\par :     1 有 WAL gy\par OM the   iM 4 dy +      5 4         L.1y)   A”        ,\par al             \textasciitilde{}  4   ITLLLE5六和气\par (fa Se ap ED (5-389) F    |\par 20 Yo          s                         S33, AS\par .         和                                    e7" iw@\par .                             萌喝线上快印\par \par 【第 18 页，来源：ocr】\par ESRRAREITA] —\par , IRA alIO SI OSCCOMmEE\par 和\par “Sy NG@\par 四    ”戎喝线上快印-
\end{enumerate}
\chapter{第9讲：静定结构位移计算}
\begin{tcolorbox}[title=解析总览]
位移计算以虚功原理为核心，常用单位荷载法、图乘法或积分法。
\end{tcolorbox}
\section{作业及答案：结构力学基础与技巧班第9次作业及答案（静定结构位移计算）}
\begin{enumerate}[label=\textbf{\arabic*.},leftmargin=2em]
\item \TODO{OCR 自动转写，结构图/公式需对照原 PDF 校对。}\par 【第 1 页，来源：ocr】\par 六、练习题
\item \TODO{OCR 自动转写，结构图/公式需对照原 PDF 校对。}\par 4-1 图4-43 所示结构中截面A 转角 〈设顺时针为正) 为(  ， )。(了哈尔滨工业大学\par 2009)\par C. 5Fpa?/(4ED                D. —5Fpa?/(4ED
\item \TODO{OCR 自动转写，结构图/公式需对照原 PDF 校对。}\par 4-2 图4-44所示悬臂粱，妃点的竖向位移为      。 (合肥工业大学 2010)\par 9   Fp=ql\par gow |                7     本      3              ，\par |一”| ”|                         a                    Ovi\par Sa,\par Bl 4-43                  .        Bl 4-44           > Do
\item \TODO{OCR 自动转写，结构图/公式需对照原 PDF 校对。}\par 4-3 图4-45所示梅架，C 点的水平位移〈 ). RRA KF 2008)      Ey ;    eas\par roe eS @\par A. 向左                   B, 向右                      FLAN\par 152                                                              BA eee FIRED\par \par 【第 2 页，来源：ocr】\par SL ee rare TA) fe\par RRS OE FES TSE\par C. 等于零    D. 不定，取决于4,/A; 的值 -
\item \TODO{OCR 自动转写，结构图/公式需对照原 PDF 校对。}\par 4-4 图4-46所示刚架中 C、PD 两点的相对线位移等于       ，两点距离\par ALTO KR 2010)           7
\item \TODO{OCR 自动转写，结构图/公式需对照原 PDF 校对。}\par 4-5 图4-47所示梅架，各杆拉伸刚度为 EA4，DC 杆的转角为 。 (天津大学\par 2010)\par 9\par El       \textasciitilde{}\par EA,   Fo      c| 三         El,  >                      Fp\par S                                 Fo .2      SE\par 8                                          A                B\par EAy                      C\par |      1    |     |  a  |  a    a \}  a  \} |\par 图4- 45              图4-46           4-47
\item \TODO{OCR 自动转写，结构图/公式需对照原 PDF 校对。}\par 4-6 RA 4-48 Baa AL BARRE, EI=SRRM. CPR KY 2012)
\item \TODO{OCR 自动转写，结构图/公式需对照原 PDF 校对。}\par 4-7 画出图4- 49 所示结构的弯矩图，并求 C 点的水平位移。设各杆 ET 相同，且忽\par 略其轴向变形。(湖南大学 2012)  i   .
\item \TODO{OCR 自动转写，结构图/公式需对照原 PDF 校对。}\par 4-8 试求图4- 50 所示结构中截面 B 左右两侧的相对转角。ET一常数。 (福州大学\par 2010)\par Cc     D\par Fp Ft Fp                  r9     O\par \textasciitilde{}                        s         1OKN/m      20kN\par i     0                人       B\par 1       1              a ©         了        |  4m  |  4m  \} 2m \}\par 图4-48                  图4-49                 图4-50
\item \TODO{OCR 自动转写，结构图/公式需对照原 PDF 校对。}\par 4-9 已知图4- 51 所示析架各杆温度上升 AcC ，材料的线膨胀系数为 c。试求 天 点的\par 坚向位移。(福州大学 2012)                               re
\item \TODO{OCR 自动转写，结构图/公式需对照原 PDF 校对。}\par 4-10 计算图4- 52 所示结构中A、B 两点相对水平线位移，ET王常数。(华南理工大\par 学 2012)
\item \TODO{OCR 自动转写，结构图/公式需对照原 PDF 校对。}\par 4-11 计算图4- 53 所示和荷载作用下的静定柏架。 (1) RAMA; (2) 求结构C 点\par 的水平位移 Ac和坚向位移Av。已知: SHRI, BD EA一常数。(长安大学 2008)\par q   s            Fe       Fp\par K                  ga   a                        E      刀   F             Mae 1 1 7\par s               8                      .O ar  Oo,\par 4              B                         Cc:       B        v2 =         二一-\par .   a  a              a    a           |      3mx4      |    yes ist\par ae      .            \_         “27. AN “\textasciitilde{}\par «4-51           Rl 4-52             图4-53        oF JO\par 1538 see FIRED\par \par 【第 3 页，来源：ocr】\par af 研究生入学考试辅导丛书 结构力学〈第二版)         CAL CNGOOSNE
\item \TODO{OCR 自动转写，结构图/公式需对照原 PDF 校对。}\par 4-12 图4-54所示结构由于支座 A 沉降 lem 引起截面 CNHRAOC=\_\_\_\_.. (B\par 庆大学 2012)
\item \TODO{OCR 自动转写，结构图/公式需对照原 PDF 校对。}\par 4-13 图4-55所示多跨静定梁支座 A 发生线位移<王10mm，转角位移 p一0. 001rad,\par 图中 l=10m, WE B 两侧截面的相对转角为 。(中南大学 2009)\par B  Cc                                 .\par A                        s             aaa\par ten] 7”                           > foe\par |   4m  | 2m | 2m |                |    7   |    1   |\par 图4-54                        图4-55
\item \TODO{OCR 自动转写，结构图/公式需对照原 PDF 校对。}\par 4-14 计算图4- 56 所示结构由于支座移动产生的铵 C 的角位移。(湖南大学 2008)
\item \TODO{OCR 自动转写，结构图/公式需对照原 PDF 校对。}\par 4-15 求图4- 57 所示结构A 点竖向位移Aw，匹[一常数。(西南交通大学 2007)\par C                       4\par A          B               \textasciitilde{}) 4\par D      7a\par k=ELIP\par 4m    om 0.01m                   I       1\par 图4-56                        图4-57\par 七、练习题答案
\item \TODO{OCR 自动转写，结构图/公式需对照原 PDF 校对。}\par 4-1 C,
\item \TODO{OCR 自动转写，结构图/公式需对照原 PDF 校对。}\par 4-2 列F (Cy)。
\item \TODO{OCR 自动转写，结构图/公式需对照原 PDF 校对。}\par 4-3 D,\par \_, gt,
\item \TODO{OCR 自动转写，结构图/公式需对照原 PDF 校对。}\par 4-4 24E1; 增大。\par (2二V2)Fr
\item \TODO{OCR 自动转写，结构图/公式需对照原 PDF 校对。}\par 4-5    SEA   ().
\item \TODO{OCR 自动转写，结构图/公式需对照原 PDF 校对。}\par 4-6 Am一中时 (<—),
\item \TODO{OCR 自动转写，结构图/公式需对照原 PDF 校对。}\par 4-7 村矩图如图4- 58 HER, Avo VIE (+),\par \_。 320
\item \TODO{OCR 自动转写，结构图/公式需对照原 PDF 校对。}\par 4-8 FF QO.                                 lay
\item \TODO{OCR 自动转写，结构图/公式需对照原 PDF 校对。}\par 4-9 一at (个)。            .                                ON NZ,\par 91ga!                                             W225    Zo.
\item \TODO{OCR 自动转写，结构图/公式需对照原 PDF 校对。}\par 4-10 24ET “7s                                  ao Eo
\item \TODO{OCR 自动转写，结构图/公式需对照原 PDF 校对。}\par 4-11 各杆轴力如图4- 59 BFA. Anc=ZEE (>); Ave=0  (BSB BIE gD\par 154                                                     BA eee FIRED\par \par 【第 4 页，来源：ocr】\par DP Be Hh Bh fie a it tn\par are gale Sill TSS ea oxen ea OE\par 特点判断) 。\par C Fycp=qa\par qa 站        OD\par U/,\par L.5qa    po   qa                       oO 9\par 1.5qa?                            “FA 了 [人wh\par 5q\par A — | B    E é           O—O          e          四\par ey  1.5qa’  Ss              四               p     37      iF,    多\par AM图                                   Fyp\par Al 4-58                                   Kl 4-59
\item \TODO{OCR 自动转写，结构图/公式需对照原 PDF 校对。}\par 4-12 0.005rad (),
\item \TODO{OCR 自动转写，结构图/公式需对照原 PDF 校对。}\par 4-13 0.003rad (4).
\item \TODO{OCR 自动转写，结构图/公式需对照原 PDF 校对。}\par 4-14 gc=0.003 3rad ()).\par \_        — 3ql"\par 4   15   Ava   SEI  ( + de\par a)\par Ne    a\}\par =A TRY £0)\par o, bare .\par OTN)\par 155戎锡线上快印
\end{enumerate}
\chapter{第10讲：力法（一）}
\begin{tcolorbox}[title=解析总览]
力法先选赘余力，建立基本体系和柔度方程，再叠加内力。
\end{tcolorbox}
\section{作业及答案：结构力学基础与技巧班第10次作业及答案（力法1）}
\begin{enumerate}[label=\textbf{\arabic*.},leftmargin=2em]
\item \TODO{OCR 自动转写，结构图/公式需对照原 PDF 校对。}\par 【第 1 页，来源：ocr】\par af REAPS REE BAA a 全 LDPECANCOA\par 1. 基本概念\par S-1 判断题。力法方程，实际上是力的平衡方程。( ) (福州大学 2012) ，
\item \TODO{OCR 自动转写，结构图/公式需对照原 PDF 校对。}\par 5-2 判断题。图 5- 125 (a) 所示柏架结构可选用图 5 - 125 (b) 所示的体系作为力\par 法基本体系。( (西南交通大学 2008)
\item \TODO{OCR 自动转写，结构图/公式需对照原 PDF 校对。}\par 5-3 图5-126 (b) 所示结构可作为图 5 - 126 (a) 所示结构的基本体系。(  )\par (烟台大学 2013)\par Fp               Fp             |                x\par 本               |                              s         ‘e                 .\par Fp      .     Foo    二一 ”|\par (a).              , \&)              (a)               (b)\par .         Fl 5 - 125                          Fl 5 - 126
\item \TODO{OCR 自动转写，结构图/公式需对照原 PDF 校对。}\par 5-4 对比图5-127 (a), (b) 两个刚架的关系是 〈， )。(中南大学 2010)\par A. 内力相同，变形也相同         B. 内力相同， 变形不同\par C. 内力不同，变形相同          D. 内力不同，变形也不同
\item \TODO{OCR 自动转写，结构图/公式需对照原 PDF 校对。}\par 5-5 图5-128所示结构的超静定次数为〈 ). CPR K 2009)\par A. 7          B. 8     Ce          D. 10”\par q               q       :\par 2EI    4EI                       .      .\par El   .   EI       2EI      2EI  i  本\par ，                          图5-127 -                  ， 5-128
\item \TODO{OCR 自动转写，结构图/公式需对照原 PDF 校对。}\par 5-6 图5-129 (b) 为图 5 129 (a) 所示结构的力法基本体系，各杆 EIS WR, b\par 为弹簧刚度，则其力法方程的系数、自由项和右端项分别为: 9 一，Ap=\_ ，\par A=\_\_\_. (HIEK 2009)                         .                     、、41;/\par 5 -7 判断题。图 5 - 130 所示结构中，梁 AB 的截面 ET 为常数，各链和的 EA 相同XGN 人\par 当 ELAR, MURR D 弯矩代数值Mp 增大。( . ) (中国矿业大学 2011)      7 Lot
\item \TODO{OCR 自动转写，结构图/公式需对照原 PDF 校对。}\par 5-8 图5-131 所示两结构中，荷载 Fe (FATE A, Bab, m. m 均为比例常数:2     int\par .      ¢          pee               ,            ”ZKA NAP \textasciitilde{}\par WH m>m 时，跨中弯矩〈( 。 )。(浙江大学 2010)                          O78\par 228     |                                                   萌锡线上快印\par \par 【第 2 页，来源：ocr】\par LOD门人二CC的下 una lip\par Fp                     Fp                           .\par 、       \_                      Fp     Fp\par k                                    I       y     “x\par |   2a |   2a    |                 x  .\par @                     )\par 图5-129                          图5- 130\par S-9  结构上没有荷载就没有内力， 这个结论在什么情况下适用? 在什么情况下不适\par FA? 〈四川大学 2009)
\item \TODO{OCR 自动转写，结构图/公式需对照原 PDF 校对。}\par 5-10 判断题。图 5 - 132 所示两匀拱，已知支座 也处下沉了As一0.01A，此时的水平\par 反力为零，尼[一常数。( ) (天津大学 2011)                    本
\item \TODO{OCR 自动转写，结构图/公式需对照原 PDF 校对。}\par 5-11 图5-133 所示结构，为了减小基础所受弯矩，可 (   )。 (西南交通大学\par 2012)                        .\par A. piv) I,       B. HK I,       C. 增大 I,      D. Ly 1, WR) 1 fF\par .   Fp   ss      ,   Fp    ,\par A      ,        B                                    q\par EI              EI\par ;                        “|           C                h\par nEl    mEl  we    ‘mE    mEI   se\par w\par A              B    h      h\par 图5-131                  图5-132           图5-133\par 2荷载作用下的计算      本
\item \TODO{OCR 自动转写，结构图/公式需对照原 PDF 校对。}\par 5-12 用力法计算并作图 5 - 134 所示结构的 MA. EI=\#\%M, EA=6EI/2, (北京\par 工业大学 2011)
\item \TODO{OCR 自动转写，结构图/公式需对照原 PDF 校对。}\par 5-13 用力法计算图 5 - 135 所示结构，并作其弯矩图，设各杆 BT 相同。 (湖南大学\par 2012) .    .\par 5- 14 用力法作图 5 - 136 所示结构 M 图。(华南理工大学 2012)
\item \TODO{OCR 自动转写，结构图/公式需对照原 PDF 校对。}\par 5-15 用力法计算图 5 - 137 所示结构，并作 M图。各杆 EI= BR. WALA 2010,\par 烟台大学 2013)     Cette
\item \TODO{OCR 自动转写，结构图/公式需对照原 PDF 校对。}\par 5-16 用力法求图 5 - 138 所示梅架支座 B 的反力。各杆 EA 相同。 (山东建筑大学 ONLI\par 2008)                      4                                   TL 从 Som
\item \TODO{OCR 自动转写，结构图/公式需对照原 PDF 校对。}\par 5-17 图5-139 (a) 所示结构，力法计算时采用如图5- 139 (b) 4 ao\par ° 3 fe. un  只\par 作其 M图。(中国矿业大学 2010)                                    oF [VO\par 22戎和锡线上快印\par \par 【第 3 页，来源：ocr】\par “il FRENESRREAR 结构力学〈第二版)             > AG mANCNar IOC\par 2kN/m         3kN\par EA\} \textasciitilde{}\par Fp      .                10kN     8         ;        i\par qr4kNim          °\par EI                                      al.     F    .\par EI       \textasciitilde{}\par \&\par |u|   1   |  4m  | 2m | 2m |        |    6m    |. 2m   :\par 图5-134             图5-135              图5-136\par Mo                                   10kN\par EI\par 8    D     E\par :   g   EI     Ey\} \&\par g\par 3m   3m  im           °\par a     \_—          @             四\par 图5-137          图5-138                图5-139\par 3. 对称性的利用
\item \TODO{OCR 自动转写，结构图/公式需对照原 PDF 校对。}\par 5-18 图5- 140 所示结构中各杆 已一常数，在给定荷载作用下，RFa = >» Va=\_\_».\par Mam=\_ 。(浙江大学 2009)
\item \TODO{OCR 自动转写，结构图/公式需对照原 PDF 校对。}\par 5-19 图5-141 aie Pee Pu =    > Fr=    。 《西南交通大学\par 2010)
\item \TODO{OCR 自动转写，结构图/公式需对照原 PDF 校对。}\par 5-20 已知图5- 142 所示对称结构中AD杆A 端的弯矩 Map = G2? /36 (AMA,\par 轴力 Fra =—5ql/12, Wi) Me=- 。(下侧受拉为正) (浙江大学 2010)\par Fp\par C             9\par 一         10kN   PTT TTT Ty\par D        人      2    |           A  B       Cc\par a                                \&\par ;                                                           D          E          F     AN LZ\par tale oly                       TFTHTTTI SN\par 4   Yy         多                                     q          \textasciitilde{}.      Car\par Va  !   u     | "2m    2m  |    a ‘e =二:-\par 图 5- 140             图5-141               图5-142    “Cine\par 230                                                   fA 228 FRED\par \par 【第 4 页，来源：ocr】\par LDL  人
\item \TODO{OCR 自动转写，结构图/公式需对照原 PDF 校对。}\par 5-21 判断题。图 5 - 143 所示对称结构 已IT一常数，中点截面 C 及AB 杆内力应满足\par M=0, Fo40, Fy=0, Fras=0. ( ) (西南交通大学 2010)
\item \TODO{OCR 自动转写，结构图/公式需对照原 PDF 校对。}\par 5-22 用力法求图 5 - 144 Prax Bitap CD 杆的轴力 ，并作受弯杆件的弯矩图。设受弯\par 各杆 了IT一常数，而 CD 杆的EA王co。(北京交通大学 2011) -\textasciitilde{}
\item \TODO{OCR 自动转写，结构图/公式需对照原 PDF 校对。}\par 5-23 用力法计算图 5 - 145 maw, HAMA. BH EA=EI/16, El=¥R.\par (合肥工业大学 2008)\par 2 二|                    人\par E      G     F\par \& .\par | |              mm            PC             DI                :                            §\par 到              |                               \&|\par me      EA                                 |       ，\par my          Bl | q         R4        可      aan      7\par J       -     Tam | am |     [tm “| am | am 2m\par 图 5- 143              图5-144                    图5-145
\item \TODO{OCR 自动转写，结构图/公式需对照原 PDF 校对。}\par 5-24 图5-146 ima, EI\textasciitilde{}\#R, FAKE MA (利用对称性) 。 (河海大学\par 2011)                              .
\item \TODO{OCR 自动转写，结构图/公式需对照原 PDF 校对。}\par 5-25 用力法作图 5 - 147 所示结构 MAL, SAF ET 相同。(河海大学 2009)
\item \TODO{OCR 自动转写，结构图/公式需对照原 PDF 校对。}\par 5-26 用力法计算图 5 - 148 所示结构并作弯矩图。(西安建筑科技大学 2010)\par 4                           10kN/m         EEN\par A        G      B\par 8              5             a        56KN B          D        F\par B                BO             EI          EI       ;\par c        F      Dd                                             EI              2EI    El     f=\}\par EI               s                            §                                      °\par E                  A        D       F                   A           C        E\par |     a     |      a    |             |    4m   |    4m    | 3m     |     3m         .\par 图5-146                图5-147                  图5-148
\item \TODO{OCR 自动转写，结构图/公式需对照原 PDF 校对。}\par 5-27 用力法计算图 5 - 149 所示结构，并作 M 图。各杆 有IT一常数，EA=ETI/2。(广-\par 州大学 2007)
\item \TODO{OCR 自动转写，结构图/公式需对照原 PDF 校对。}\par 5-28 利用结构的对称性作图 5 - 150 所示结构的弯失图。已知 已T一常数，链杆EA为 | AN 1 '\par 无穷大。(中国矿业大学 2009)                                                                   ONLI
\item \TODO{OCR 自动转写，结构图/公式需对照原 PDF 校对。}\par 5-29 用力法作图 5 - 151 所示结构的弯矩图。ET一常数。(了哈尔滨工业大学 2013) Te!  Lo.
\item \TODO{OCR 自动转写，结构图/公式需对照原 PDF 校对。}\par 5-30 用力法分析图 5 - 152 RRM, AMA, RSSBH ET 值相同。  (清华大学      cao\par 2                                          .                            ts. Ay  .\par 23戎和锡线上快印\par \par 【第 5 页，来源：pdf-text】\par 叫研究生入学考试辅导丛书结构力学(第二版)\par EA\par C\par q勹B\par G§1』\par 20kN\par IF\par 丫G\par y\par A\par 7?7?77' 玉\par j\_\par 2\par 尸\par B\par ＼\par L\par L\par 6m\par 6m\par l\par 1\par 1\par ．，\par 图5-149\par 图5-150\par 4. 弹性支承超静定结构的计算\par \par 【第 6 页，来源：ocr】\par EBA R Stil\par RO TUBTODIGHVSSS\par 52. V
\item \TODO{OCR 自动转写，结构图/公式需对照原 PDF 校对。}\par 5-3, X           C\par 8\par » A      age \#5, 014 Bs\par A             aw AG  pr”\par be\par sx.\par Spe\par \par 【第 7 页，来源：ocr】\par EEL SBAR STN\par ISM OB IGBUSSS
\item \TODO{OCR 自动转写，结构图/公式需对照原 PDF 校对。}\par 5-5,\par 和\par Hy    fils Ate,\par ,       |             bn= Fen opel =\par yA in nl >\par 4:             be\par 5, CA, A RAY pRB adi,\par ny N\textbackslash{}p K )  x\par 3s bots, fae     » Mae Me,\par 45197    i\par My \_ we\par A RAAT RO po A\par Ne 中 bb \& Her SF).\par on\par op        CL yi        \%.\par ‘yt         an      Lad ae BT\par CN)\par \par 【第 8 页，来源：ocr】\par HELSEAR aT\par MCAvin aN O31O37666\par tll, Lo af AMSy 博大II , Red,\par ph a it\par b-42,\par L         -|    ig            ‘            =   +L |\par a\par ips — el Lx *).-  fa          :           和\par duit — 4  S  4 xi * De a | 本\par a  ft\par ff       4\par \%\par T\par was\par SO\par \par 【第 9 页，来源：ocr】\par HE SEAR Sis\par ISM OB IGBUSSS\par ik\par orn\par \textbackslash{} |    ?      8               本\par h    m1          we? fur\par bic if  Ne xy ]+ Erp \textasciitilde{}\textasciitilde{} 本\par \textasciitilde{}      1站\par be Lt Blt Qe =\par - By [gare Bt 497]\par bipe - 218\par 5\par bu UF ep   47 3\par thy\par ve\par 13.06    |\par re] 2  |\par m mp            ,\par 4141.7\par we\par SND)\par 第175页, \$359                                         meade ERED\par \par 【第 10 页，来源：ocr】\par RAR Stil\par MAiMaiosO3sTeSs\par Say,               本\par ay\par ih           :      、 El\par 6        ‘          il     Eu      Lad “F\par 1\par ae |                    =\par ;          a CE\par Pipe A" 如     - Fy ¥ 9\% br (bor) ¥ =  y\par Suxit Spo    eye 346\par abAY¥\par b\par yb   7.16\par :\par sy, ® \%          as\par of\par ar,               Mp\par diz AE Astin 4] =  a\par Sipe 414\% 及站 学 一 th A+ + ms村| 二一 oe\par a)\par ONL\par OPO\par \par 【第 11 页，来源：ocr】\par EAR Sis\par ROGUE OD IGBIBSS\par pivot 21p29  a 为= ae . Ls \_ =\par Ano   AN\par ¥                | J Mo\par Me\par +\par tab,                 adi\par |< bet          >               ,\par 54    seat\par ra   L oR F\par x71  山               Th\par mi              Np\par kar]  dnt 24\par bys th [ay y+ Sry DAF 7| ioe:    \_\par Ape 此     va -和At \textasciitilde{} Sox Deas fF ]2 - We\par > wl -h\par Gy) hE - », is\par win QE \%  p, bstfp  Ls\par 54],                   4\par 4"\par irs       ne\par mM            mp fos.\par v         =\par 四\par + abit wed] 页\par 11 .7/\par we\par OPO\par 第177页, F£359                                         asad ERED\par \par 【第 12 页，来源：ocr】\par 后博美技公众强大]\par 未久微信19818587888\par 4     >   ‘\par +          a 徊> 学\par Le        |\par 区\par iv   ei      |\par 7 4 国7                  站\par a+  + We Bde +) -Fel ‘\par Ape Al -      - tp- as] = sss =\par 4 Lpe be tee P24 Sess (34 94)\par 7 aby 4x 147 + Px bes 二|\par ”   o        4s+ as|= =\par ye el At ‘ot           区\par 本\par AL二 下[ly ax lt 444) \_trby “ ,*\par VSP ty4)] =\par acl x    2   ogee 中\par Ls 六                Yaz — 4s.\par 3\par if WW\par Ye\par a\par y            |\par pa °G\par 11.7\par AND\par \par 【第 13 页，来源：ocr】\par SEAR StS\par TRAE OS IOSESS\par bai, aad fh\par bi. fT      NU =\par fay x Le Re 2)\par sys Xf\par "ete 5      Hac - fr     Vz tr\par 544.        le          44.\par 5 yt  +H\par MAO\par 4 Kk     Ayxze SHe Ss pies\par o' ¢            1
\item \TODO{OCR 自动转写，结构图/公式需对照原 PDF 校对。}\par 1-1 1h 4\par Sw ;     ia Mya Maw v4 Yxe 42\par a. Meg,\par ye Vii, BAY wy         |\par ya ee\par i\par "    ¥     5\par bie + + Ws -\par we\par bal, V\par NO\par \par 【第 14 页，来源：ocr】\par HELSEAR aT\par MCAvin aN O31O37666\par ty\par wv,              \%\par Hy\par IK I<       4)         "\par ;                                       |\par Y,           人\par ”»                wy     she  +\par SiPie 2\par 和PP>= aly   et lo al\par Wy            Sipe  所 |\par Bays\par 如 出 ls    MT ayo  \& Y=  \textasciitilde{} *k\par a7\par a       y\par y), A 4  ak\par Xr\par s                 be\par wlan.\par ONL\par GRAND)\par \par 【第 15 页，来源：ocr】\par HEL RAR aii\par TRAN SILO OSTOCS\par —\_—                               ，\par 2\par +I\par ott\par 1                    Wo\par ;                 i\par L                  pm\par hie i.   aw le re 4 |\par 1    by     3\par we GL ee al> 抽\par L          Lo +]       pe\par bye 1h [ s* wor2*F Yt    a7 ma\par a OA\par bipe al +" 2HF wit |- x\par bye -hl \$x noe | =一  总\par 32        ib       \_ 虹            44\par a)     bg      "   " aA\par -    Yt     Z   tL              9\par 个            ¥¥\par ww Baws\par 1h             | 由\par AN    ab\par 11 .7\par NO\par \par 【第 16 页，来源：ocr】\par EAR Sis\par ROGUE OD IGBIBSS\par fof,                     \_            \_\par Akhiok      相生 各           党\par Ale\par a>        i\par basse          \textbackslash{}\par 2”      x4                       2\par den           x                  oe\par BD                            a            4M,\par be Blt ay 44让- SE\par lex 442< bo boom              1                 .              ;\par ¥ 4° 0 kun  bya alte hx MInZ+ ty qu了]\par r—4,          [下\par se                 i=    s61           a\par i     .           t + USxwe »    \_4  :\par = 4\par ?人\par owe GL ae bebe dey - Fe MHD I GL ee — OO). Se\par wlan.\par SND)\par \par 【第 17 页，来源：ocr】\par EBA R Stil\par 未久向信1981387888\par Ad\par BS, is we x,  - A      |\par et          wel          17 0\par | tn  fae OP\par AN        am\par 5-H,\par | i=  ii      ;\par y 于           we         It\par Sey          its\par |                                   ba ，\par .               6, ceox BRE get BBY >\par .                                           -经 Eady \$441-\par |            *                    SuXi4 220  2X7)\par a\par ,         i\par 有           Iz         人\par a”\par 4                     up\par | hl [pew\par wlan.\par ONE fo\par GRAND)\par Bite ASK\par \par 【第 18 页，来源：ocr】\par EAR Sis\par MAiMaiosO3sTeSs\par FU). VA hiv +49\par 2 a     i. a\par M4      i     L       L\par >  -       a\par |          .             Se 5 x Es Bee ot.\par Me Gert Me ae\par r     ,\par \_ Soc一下[二dk中:- a\par Ap Gy (tale req] ae 7\par ons\par y       L         中\par +. SSL       六 00/3\par 》       Y\par = Xi +  i)   0        Xr= 9-1)\par chew 3 条 io w\par \textbackslash{}  0.4   fw\par “\par wlan.\par ONL\par OPO\par \par 【第 19 页，来源：ocr】\par HEL RAR aii\par TRAMs Os IO3TSCS\par Sm. Rath vy        |\par HA hf. Mo\par ga\{ At , Aver, 4h   be\par >             > mm      一       fx\par RE               ht    to\par b         |\par 和\par Np,\par |           My                     jh         +       7\par |     ane Hf de babebeg t Peel) = G5 [n+ Mile a\par 本\par |     bu xf 6° P=0    > 为  Es                   日\par a                                                    |\par “7 ly /\par Pi            i)\par 11 .7\par we\par CO\par \par 【第 20 页，来源：ocr】\par HELSEAR aT\par 未久向信1981687888\par SA. jetted The 41K\par k-          并\par 了 |        |  ll\par w= -4       7\par 2\par wy\par ta.\par tt iow tt le iow        人们mw\par ei            i\par “\textasciitilde{}\% y P.-.                 =\par UF\par uke\par =F                                 IX\par ee ao il\par re         1S     7\par i      Zo.\par mt\par ES           :           ur\par bee - by [ Liye lS plOKy + 8 24204 tk]2- A\par byt Spay eh\par 4141.7\par RN\par \par 【第 21 页，来源：ocr】\par ESA Rail\par RORY OBIGSTB88\par b    . b\par 6\par b                      ‘\par \textbackslash{}v ie\par ¥   Me   \_\par RN\par Bit:RE ize
\end{enumerate}
\chapter{第11讲：力法（二）}
\begin{tcolorbox}[title=解析总览]
力法（二）重点检查柔度系数、自由项符号和变形协调条件。
\end{tcolorbox}
\section{作业及答案：结构力学基础与技巧班第11次作业及答案（力法2）}
\begin{enumerate}[label=\textbf{\arabic*.},leftmargin=2em]
\item \TODO{OCR 自动转写，结构图/公式需对照原 PDF 校对。}\par 【第 1 页，来源：pdf-text】\par 叫研究生入学考试辅导丛书结构力学(第二版)\par EA\par C\par q勹B\par G§1』\par 20kN\par IF\par 丫G\par y\par A\par 7?7?77' 玉\par j\_\par 2\par 尸\par B\par ＼\par L\par L\par 6m\par 6m\par l\par 1\par 1\par ．，\par 图5-149\par 图5-150\par 4. 弹性支承超静定结构的计算
\item \TODO{OCR 自动转写，结构图/公式需对照原 PDF 校对。}\par 5-31 图5-153 (a)所示结构，取图5 - 153 (b)为力法基本体系，其力法典型方程\par 为(\par )。k为弹簧刚度。(中南大学2012)\par Fp\par a\par 十\par 图5-151\par a\par 一。\par 一\par 20kN\par I\par EA==\par \{\par 3m\par \{\par l\par r\par k-\par (a)\par (b)\par 图5-152\par 图5-153
\item \TODO{OCR 自动转写，结构图/公式需对照原 PDF 校对。}\par 5-32 求图5-154所示结构在给定基本体系下的力法方程及方程中的所有系数。(青\par 岛理工大学2008)
\item \TODO{OCR 自动转写，结构图/公式需对照原 PDF 校对。}\par 5-33试计算图5-155所示具有弹性支座的结构，并作弯矩图。抗转刚度系数ko =\par 單(中国矿业大学2008)\par q\par q\par 贮I I I\par I I I I\par 贮A\par u\par j\par C\par EI\par EI\par I\par ］\par ko\par M\par El嘴数\par k\par \{\par lr-3El/13\par ＼\par 1-\par f\par I\par 112\par I\par //2\par 图5-154\par 图5-155\par 5. 支座发生位移和温度变化时超静定结构的计算
\item \TODO{OCR 自动转写，结构图/公式需对照原 PDF 校对。}\par 5-34 图5-156中等EI值单跨梁两侧温度变化时，M图位千(\par )。(青岛理工大\par 232\par \par 【第 2 页，来源：ocr】\par Mt a Be sh Sh\par AY) rnin\par A. 低温侧       B. 高温侧       C. 不同侧       D, =\par 、        mle     一                       Ou X1 十92十Ac一A
\item \TODO{OCR 自动转写，结构图/公式需对照原 PDF 校对。}\par 5-35 判断题。 FLAS  157 所示结构基本体系， mA X,+0nX,-+Ae= 人中，\par Aic=0          ve                ，\par ae C  ) (福州六学 2011)
\item \TODO{OCR 自动转写，结构图/公式需对照原 PDF 校对。}\par 5-36 图5-158 所示刚架除承受荷载外，支座A 下沉了4二2cm，且顺时针转动了0一\par 0. 0lrad。已知 E=2.1X10'kN/cm?, [=2000cm', FAA ERICMIZER SA. 华南理工\par 大学 2009)                              ，\par 加       oy        g\par <              |      50KN              “\par Xx\par -                                   和\par +2rC                                              A\par 上一一一一才           (a)                )          上一全\par 图5-156                    图5-157        ;   Py 5-158\par 6. 超静定结构的位移计算
\item \TODO{OCR 自动转写，结构图/公式需对照原 PDF 校对。}\par 5-37 图5-159 所示梁 BK) 支座位移产生的跨中点竖向位移为(  ). (aK\par A. lp/16        B. lp/8 - OG. Up[人4        D. lp/2
\item \TODO{OCR 自动转写，结构图/公式需对照原 PDF 校对。}\par 5-38 图5-160 Mani Re HRA EWRHSHA, WC AHAMB oc=\_\_\par \_\_ ，五点的水平位移Anz一        。 (湖南大学 2011)
\item \TODO{OCR 自动转写，结构图/公式需对照原 PDF 校对。}\par 5-39 5-161 所示一端固定、另一端定向支承的等截面粱，当 A 端发生单位转角\par 0一1时，B 端的竖向位移A一           » (BRK 2010)     -\par 11.98\par 7.29      Fp    3.65                  ,\par El            EI  E          .\par ，                            B       c     8,33    1.04\par EI      2EI         EI    \textasciitilde{}     ot          .\par ¥     EI Ph\par A | \textbackslash{}365     D.  ” Fih\textbackslash{}o.s2      ee         Fl\par 9    EI            a     ，           “oa BS -  、 0\par x eae                             MPA(«0.01F pl)          上 -|            \textbackslash{}QN: 2\par 图5-159                 图5-160                       图5- 161        cht 二\par 5- 40 图 5 - 162 所示结构，各杆抗村刚度 ET 为常数，忽略杆件的轴向变形。(1) RN\par e  ; 1 ' .\par 23稍揭线上快印\par \par 【第 3 页，来源：ocr】\par app RS AE RAR BRAS Aa a LEP eaASE\par 据结构与荷载的特点，选用最简便的方法作刚架的弯矩简; (2) RC 点的水平位移。 (武汉\par 理工大学 2011)            ，                ，\par 10kN\par ，               20kN-m            .\par |                C     D      E        F\par \&\par ’               A             Bl:                   .\par 2m   2m   |   2m\par 图5- 162\par >\par 234                                                          fA eee FRED\par \par 【第 4 页，来源：ocr】\par RAR Stil\par 未久向信1981687888\par SAI. buryt OnE\par Sd.           \textbackslash{}   |\par 人              yi\par bs                    加\par et                    Ay\par i             Sue Galt (wy eet cs\par yw          |\par be aia + Oe a\par ®            Ine Sl dnvetrlag|+ 这\par > of\par AN       Ine 5      = 7"\par £.\par > 3》\par 外 aN  ee    和   BL - th\par aes   .    》       a      5     \textbackslash{}\par a i dyldele Fo bt J     a ae a? - eye\par bat ya a 4\par bart SXrt 54 = 0°\par a)\par OPO\par \par 【第 5 页，来源：ocr】\par HELSEAR aT\par MAiMaiosO3sTeSs\par ae\par Pr                ‘i cate\par bye 4 \%   MN bale nz] 2 al          =\par se, ih - acd ,  x My 7 stn tay My   +    Ll   xF\par suit a oa  ts YN 0 ek\par vv Ye    A\textbackslash{}\par oy, A            oe      .\par Fr,                    Fx    从\par |           -                             .\par L   人  Sun                  4=-b.\par y                 ‘ol\par a)\par ONL\par OPO\par \par 【第 6 页，来源：ocr】\par HELSEAR aT\par IRB IESTSSS\par ri                         yuz 4 [ax douse 5 t x att]\par 4        Tw                         ，\par 4               局- 10.37 mw\par TM te aa\par wont 儿    any v| 9  s bes \_pabm\par \$y Xt     S5iL=     \textbackslash{}" by\par ped.\par Bore\par er        ia |
\item \TODO{OCR 自动转写，结构图/公式需对照原 PDF 校对。}\par 5-4]. B          rh ¢          \textbackslash{}y       1\par =   tl\par Ay\par 17\par (但 + f  2H vy diego  AM\par ate      yl\par |           f= tt x   人\par 才 tls (4+ N48) xi]- wee Lx Gry\par |                     Ye F sein\par A\par 4141.7\par we\par OPO\par \par 【第 7 页，来源：ocr】\par RAR Stil\par MAiMaiosO3sTeSs\par aa |       .\par chtc2 4, Ly 14x ley 一 jx Lx 256 bt]\par bape — Mele Rl\par L                                        aL
\item \TODO{OCR 自动转写，结构图/公式需对照原 PDF 校对。}\par 5-4,  -\par b= \textasciitilde{}    * Lebe    t) =   >\par - fo、 ()\par ph             2 my       Toke\par Ad               ry          (¥\par |         7\%\par L\par im                  hy> als mre t wer\}\par Ww\par 9            32\par *e        ie BL  + = 3h\par y                        a\par ib       yx2ax W) - \_           bnXi+ shay G  j=  =z\par » |\par if  >               a     a\par M                2          is\par QUX= alts rn] =  =\par wlan.\par ONL\par OPO\par \par 【第 8 页，来源：ocr】\par app RS AE RAR BRAS Aa a LEP eaASE\par 据结构与荷载的特点，选用最简便的方法作刚架的弯矩简; (2) RC 点的水平位移。 (武汉\par 理工大学 2011)            ，                ，\par 10kN\par ，               20kN-m            .\par |                C     D      E        F\par \&\par ’               A             Bl:                   .\par 2m   2m   |   2m\par 图5- 162\par >\par 234                                                          fA eee FRED
\end{enumerate}
\chapter{第12讲：位移法}
\begin{tcolorbox}[title=解析总览]
位移法以结点位移/转角为未知量，由杆端力和结点平衡建立方程。
\end{tcolorbox}
\section{作业及答案：结构力学基础与技巧班第12次作业及答案（位移法）}
\begin{enumerate}[label=\textbf{\arabic*.},leftmargin=2em]
\item \TODO{OCR 自动转写，结构图/公式需对照原 PDF 校对。}\par 【第 1 页，来源：ocr】\par 三 的[Re\par 十、练习题\par 1. 基本概念
\item \TODO{OCR 自动转写，结构图/公式需对照原 PDF 校对。}\par 6-1 位移法的典型方程与力法的典型方程一样，都是变形协调方程。( ) (烟台大\par 学 2014)
\item \TODO{OCR 自动转写，结构图/公式需对照原 PDF 校对。}\par 6-2 图6-117 所示结构用位移法求解时，结点未知数为  “。(天津大学 2008)
\item \TODO{OCR 自动转写，结构图/公式需对照原 PDF 校对。}\par 6-3 指出图6- 118 所示结构用位移法解题时，未知量的数目。(武汉大学 2010)
\item \TODO{OCR 自动转写，结构图/公式需对照原 PDF 校对。}\par 6-4 图6-119 所示结构，若用力法求解，其最多的未知量个数为      ; 若用位\par 移法求解，其最少的未知量个数为       。 (西安建筑科技大学 2011)\par 五天co                                         O\par EA=co\par Fp                                             ©\par Al 6-117               Al 6-118            6-119
\item \TODO{OCR 自动转写，结构图/公式需对照原 PDF 校对。}\par 6-5 图6-120 所示结构中CD、FG 杆的ET一co，其余各杆 EI 为有限值，忽略杆件\par 轴向变形。若用位移法计算，则其基本未知量数目为  )。(浙江大学 2011)\par A. 1          B. 2          Cc. 3          D. 4
\item \TODO{OCR 自动转写，结构图/公式需对照原 PDF 校对。}\par 6-6 图6-121所示结构，若采用位移法计算 (SH ET一常数，忽略轴向变形)，则\par 基本未知量数目为， 结点角位移    个，结点线位移    个; 若采用力法计算，则基本\par 未知量数目为    个。(浙江大学 2009)
\item \TODO{OCR 自动转写，结构图/公式需对照原 PDF 校对。}\par 6-7 图6-122 所示结构各杆的线刚度;相等，则位移法典型方程的 eo 一        ，\par AR22 一     。(了哈尔滨工业大学 2009)\par I\par G\par Cc    D    F     H           了      A,         A           5\par 一、         2        43\par B       A       C      四\par G     >     E     “                           \$\par 下    D       H\par A          B     A          B       |   6m   |   6m   |\par KI 6-120            A 6-121                  A 6-122\par 2. 荷载作用下的计算                                      laws
\item \TODO{OCR 自动转写，结构图/公式需对照原 PDF 校对。}\par 6-8 图 6- 123 所示结构，和欲使结点 BRO AWS, Fe ( > CAMRRR 200\%,\par A. 2qa         B. 8qa/9        C. 3qa/2        D. ga     Tr aK cay -
\item \TODO{OCR 自动转写，结构图/公式需对照原 PDF 校对。}\par 6-9 用位移法作图 6 - 124 所示结构弯矩图，忆一常数。(广州大学 2011 BYZ*  So\par EAI\par \#08                                                    萌喝线止快印\par \par 【第 2 页，来源：ocr】\par a RG A REISS (0 BIS ES VE UE G\par 这
\item \TODO{OCR 自动转写，结构图/公式需对照原 PDF 校对。}\par 6-10 用位移法作图 6 - 125 所示结构的弯矩图。(福州大学 2011)\par B           Fe       C                                           \textasciitilde{}\par eT                              q          q                      fr            。              >\par E                     2EI\par ?| ler                       8                                      \textasciitilde{}\par EI          El) |\par 4                                                       ;                                Al            B   |\par |    a     |     a     |              |     7      |      7                    |       I        |        7       |\par 图6-123                图6-124       图6-125
\item \TODO{OCR 自动转写，结构图/公式需对照原 PDF 校对。}\par 6-11 用位移法计算图 6 - 126 所示结构，作其弯矩图。各受弯杆件 ET一常数。(湖南\par 大学 2011)
\item \TODO{OCR 自动转写，结构图/公式需对照原 PDF 校对。}\par 6-12 用位移法求图 6 - 127 所示无侧移刚架的杆端弯矩 FASE. ROMS\par 件的相对线刚度:横梁为 1，坚柱为 1.5 。(中国地震局工程力学研究所 2011)\par C                  D   ,                                     son         21kKN/m                150kN-m\par q             Fp=q!       B   -                                                            有\par A                 E\par |    I    |  Z2  |  Z12    |. 2m | 2m |    8m        6m   |\par 6 - 126                                                         fl 6 - 127
\item \TODO{OCR 自动转写，结构图/公式需对照原 PDF 校对。}\par 6-13 用位移法求解图 6-128 所示结构，作 M图。FET 为常数。(青岛理工大学 2010)
\item \TODO{OCR 自动转写，结构图/公式需对照原 PDF 校对。}\par 6-14 用位移法计算图 6 - 129 所示结构并作其弯矩图。(西安建筑科技大学 2008)\par B\par 1.5EI          €                        34kN/m\par 2     EI         D    El      A                                     3EI\par 1.SEI          后              El                      EI       \&\par Cc\par |     4m      |       Sm       |                       |             8m             |\par 图6-128 6-129\par 3. 刚度无穷大杆件的计算                                                                     NA
\item \TODO{OCR 自动转写，结构图/公式需对照原 PDF 校对。}\par 6-15 图6 - 130 所示结构中截面A BHM, =               MISERY.   BION) Eo\par fA 522k ERED\par \par 【第 3 页，来源：ocr】\par af PEAS SR A RE eB Re NEARY Mferlerve
\item \TODO{OCR 自动转写，结构图/公式需对照原 PDF 校对。}\par 6-16 试用位移法计算图 6- 131 所示刚架 (EI= BRO, KHAHRSEA, KS\par 2008)\par 12kN/m\par A   EI      B\par .                           EI                 \&\par El,=00      El,=00       El,=00       Esco Fe                           c-— =>\par EI          EI          EI          EI        3                        EI           al \&§\par A                                                                 D              E\par a     |     a     |    “a     |     a     |                      |     4m     |     4m     |\par 图6- 130                                                          Al 6-131
\item \TODO{OCR 自动转写，结构图/公式需对照原 PDF 校对。}\par 6-17 FA 6-132 所示结构的M图、Fa 图及 PVA. CPR LAS 2009)
\item \TODO{OCR 自动转写，结构图/公式需对照原 PDF 校对。}\par 6-18 用位移法作图 6 - 133 所示结构的 MA. REI 为无限大，各柱相对线刚度\par 如图中所示。 (华南理工大学 2010)\par 4                     .\par Cc     EI        E\par Eco\par El     =    EI    -           a        |        |      |   i\par 1        1\par 图6- 132                                                        图6- 133\par 4. 对称性的利用
\item \TODO{OCR 自动转写，结构图/公式需对照原 PDF 校对。}\par 6-19 用位移法求图 6 - 134 所示结构的弯矩图。开IT一常数。(北京工业大学 2013)
\item \TODO{OCR 自动转写，结构图/公式需对照原 PDF 校对。}\par 6-20 用位移法作图 6 - 135 所示结构弯矩图。天T一常数。(苏州科技学院 2007、华南\par 理工大学 2006)\par 4\par q           |           5  |        ae\par 4                        \textbackslash{}\par )                                 图6- 134                                        图6- 135         Ty 24  iz.\par 308                                                                                                                        Ch in\par fA aE FIRED\par BSii-RE ARR\par \par 【第 4 页，来源：ocr】\par 型
\item \TODO{OCR 自动转写，结构图/公式需对照原 PDF 校对。}\par 6-21 作图6-136 所示结构的M图，各杆 也[一常数。(西南交通大学 2010)
\item \TODO{OCR 自动转写，结构图/公式需对照原 PDF 校对。}\par 6-22 用位移法作图 6 - 137 RAMA. SH 已I一常数。(山东建筑大学 2012)\par 4kN/m                                           -\par pett tty ttt did...\par n\par a         I          I            \&                                            ,\par A          -                NB                                    -\par |       3m      |      3m   |   2    |   We    |                :\par Al 6 - 136                                      Al 6-137
\item \TODO{OCR 自动转写，结构图/公式需对照原 PDF 校对。}\par 6-23 用位移法作图 6-138 mA MA. SH EI=BRM. (UREA 2009)
\item \TODO{OCR 自动转写，结构图/公式需对照原 PDF 校对。}\par 6-24 用位移法作图 6-139 maw MA. BRS EI=BRM. (J KS 2010 年)\par 40kN\par 2   1/2  |   1/2  |  12   2  2  |                 |     4m     |  2m  |  2m. |    4m\par 图6-138 6- 139
\item \TODO{OCR 自动转写，结构图/公式需对照原 PDF 校对。}\par 6-2S   用合适的方法计算图 6 - 140 所示对称结构，  并作出 MA. EI-\#R. CHR\par WHA 2011)                     .                              oe
\item \TODO{OCR 自动转写，结构图/公式需对照原 PDF 校对。}\par 6-26 根据对称性用位移法求解图6- 141 所示结构，  作M图。FT 为常数。(青岛理工\par 大学 2012)                                 7\par A         C         E                        12KN ”3kNm\par F       ye                               D\par 一                            一                                               \&\par B         D         F                         A         B\par 4m   "  4m   \}             | 2m | 2m |   4m   |   4m   |\par 6- 140                                     图6-141                 Na He)
\item \TODO{OCR 自动转写，结构图/公式需对照原 PDF 校对。}\par 6-27 用位移法作图 6 - 142 所示结构的 M 图。EI一常数。(烟台大学 2013) 72F GV) =\par 235. et\par 7 NO\par .                                                      fA ZR FIRED\par \par 【第 5 页，来源：ocr】\par 了2 DEO RT ARF IAM\par 5. REV HAASSAHBK
\item \TODO{OCR 自动转写，结构图/公式需对照原 PDF 校对。}\par 6-28 画出图6- 143 (a). (b) 所示结构弯矩图的轮廊，杆件抗弯刚度 ET一常数。\par (北京交通大学 2010)\par |    +fC\par | Fp\par 7          1           1                                             ,\par |          |          |           |                      (a)                       (b)\par 图6- 142                                                            图6-143
\item \TODO{OCR 自动转写，结构图/公式需对照原 PDF 校对。}\par 6-29 定性地画出图 6 - 144 所示结构在图示荷载作用下弯矩图的大致形状。 (武汉理\par 工大学 2010)\par 6. 弹性支承超静定结构的计算
\item \TODO{OCR 自动转写，结构图/公式需对照原 PDF 校对。}\par 6-30 用位移法计算图 6 - 145 所示结构并作 M 图，中间支座为弹性支座，其刚度系\par 数 上一5ET/R ，ET一常数。(合肥工业大学 2009)\par |*\par k\par 2   |   1/2          了       |\par 图6- 144 6-145\par 7. 支座位移和温度改变时的计算
\item \TODO{OCR 自动转写，结构图/公式需对照原 PDF 校对。}\par 6-31 用位移法求解图 6-146 所示结构的 M 图 〈须写出计算过程)。(已知，ET一4. 8义\par 10‘KN > m’) (东南大学 2007)
\item \TODO{OCR 自动转写，结构图/公式需对照原 PDF 校对。}\par 6-32 6-147 所示刚架，妃支座产生竖向位移A 时，试用位移法作弯和矩图。 (KE\par 大学 2008)\par 2i\par |                         4              B\par |      6m    i“      6m     |                 |        i         |    ‘\par 图6-146 6-147              Sxl Lio\par 8. OER RM                                       Bl ea
\item \TODO{OCR 自动转写，结构图/公式需对照原 PDF 校对。}\par 6-33 用位移法求解图 6 - 148 所示结构未知结点位移。(东南大学 2013) ry SO SS\par 310                                                                                                        LENO\par 萌锡线上快印\par \par 【第 6 页，来源：ocr】
\item \TODO{OCR 自动转写，结构图/公式需对照原 PDF 校对。}\par 6-34 图6-149 所示结构各杆抗弯刚度 下7 为常数，BC 杆受到均布荷载作用，同时姜\par 支座下沉A。试求 B 结点转角，并画出变形图的大致形状。(武汉理工大学 2010)\par EI\par q\par EI      Elmo\} = g       \textasciitilde{}                 7;             3              C\par 及      EI\par C £El                                                           s\par 28 | a                    \textasciitilde{}                         D\par el\par I      |      1      |             |       a      |      a       |\par 图6-148                                              图6- 149\par re\par SiR iki\par \par 【第 7 页，来源：ocr】\par ELSA Rail\par RTE OBIGSTSSS\par be. x\par h2. 17\par 4\par 63. uf\par nk\par 二\par 1 -,   \textbackslash{}   i\par +p,      :       ay   ‘BA\par i\par Fai     a”\par (4,       |    bt,     @   hb\par bh: x\par >|\par “Ge\par wire eR\par \par 【第 8 页，来源：ocr】\par EEA RS TID)\par RO TUENOD 1GBUSS8\par 3"                     q    "\textbackslash{} Pi\par 1                           方           A\par i\par > 人         中|\par Wu Wi ll\par we Ce    \textbackslash{}\par ba)                      in\par |        b Ye\par b4. ei 5 5A\par MW\textbackslash{}A>o\par 和      a”      Fox eabs vc z\par :          ates fe EG py, Be AN\par An  Tif pe 0\par 4   h     |\par wa  a         ;    di\par 一\par = ”协昌\par 和           AM\par “So\par BSi-:ReE kei\par \par 【第 9 页，来源：ocr】\par RAR Stil\par RUE ODIGSTOSS\par ou”\par Sy\par y/) We    au\par 了\par ie\par By\par \}\par ig      oh   a       42.\par Ww             ‘ql Pip Fpl         1D\par px                        2B\par We ,  ‘In + Th? *\par Te Tye 一 8.                   有\par iu\par We  Bx a> Oy    mi  37 EL\par 5\par re tbe 0"     六oz he\par |          th\par 了了    | 1\par NOD\par BRR ASST                    aasaie' LERED\par BER eS\par \par 【第 10 页，来源：ocr】\par ELSA Rail\par RTE OBIGSTSSS\par bl.      3\par que dit lt dizi\par 由         \{nz 3 Ye 咎\par Thnz> Wnz 外\par Pabz = We + :\par bre fut 押\par 10 YF YX FP         a\par |      Yn’   an dal    , |   空   5\par ,                  W\par Yor -   “Tt\par i                    4\par \% tb\par 49                 =\par el eee +                  本\par ers © ja lil\par ie       -,\par bx   > 出|\par 二 be ee\par Th       py 2 Js + sire 3] jam\par BSi-:ReE kei\par \par 【第 11 页，来源：ocr】\par ELSA Rail\par RITE OD IGBISES\par ty     2   Yer                  ¢\par 四1人2，Ynoo IS\par i 六 +   2    +     Xj= 26]\par i   WS\% Be]      2 3bR\par pl] POD ob, sp.\par app           la                       |\par ?|\par MA四[em)\par 64\}.\par Ly dy Ws   "\par \&   中          Pap Y om\par + ab\par x  My    Y\par a     oo                    bs\par 外\par wa\textbackslash{}/ 县\par 4      ae                    2B   “Fh\par a                    as\par GRO\par SiR cei\par \par 【第 12 页，来源：ocr】\par 久久极蛋1991387888\par 1a       8\par Tie —E | Tre Sp Pi: 他\par fay + By        Xz Ae\par his ol in Mintle       =\par We x+ BEB y, 2 ib    \%z \$f\par uf\par be           rr\par ey\par bHUp.\par Aw\par 1    a\par x                      BX Wee Ye\par ws Paps > lav.\par me\par @                Tne  18\par ae\par a, \%    Huet pm os Me BP Re\par *\par by\par NO\par wire eR\par \par 【第 13 页，来源：ocr】\par SEAR ais\par IAMS OsIO3TE89\par bay,\par eT TY \textbackslash{} mae by 3.8\par A\par bab,\par *|                time!\par Ne 4it die        jee\par V            uel Ths  |  v)  mi\par Hy   Yuzyye 2\par bp\par quit Yhyxz十 PP?    ,| £\%  om\par ae ‘tt Pp   ) | ue 乌\par 外\par ia\par ]\par CaN\par \par 【第 14 页，来源：ocr】\par 未久向全1981387888\par 人\par hy  \%  ow     hy +\par Bs\par +\par AN\par Vo\par Yuya Pip =0    2   hl = «seb\par we    al 证\par 也                           2 ¢€     \_— wiz\par aN”  i  ou”       at     539818 ca\par op  v        — a so\par T mB"        中        na\par bY,\par os aa i Ba —   bx ovine\par \}              Tn   =\par 4   ie       Ll\par /   Wa   /  /|   4   =\par 8\par «nyt ?    本  Ae x 20\%   \& ee   Ww?\par ree.\par je\par 2 ed ced ca\par e      we   we\par io \%   \_ im 7\par “Spo\par \par 【第 15 页，来源：ocr】\par HEL SAR STN\par RITE OD IGBISES\par 4, Hee,\par “a din                 ot t—i,\par 7       ee\par 2    b                Fi                          241\par 到       be\par ws   aL       ic\par (oe       \{wu              8\par Yu= 至   ,人和: 下+  me  .,  zf一\par a>\par a"\par 14”\par :              yi\par SY\par W | a”\par 了\par are       Ahh, 7            Pa\par ty                        a\par an\par a\par ENO\par \par 【第 16 页，来源：ocr】\par SEAR ais\par IAMS OsIO3TE89\par 62,              x" 4     3\par EX 43\%. ahem\par nde Py = 4\par FM\par ZL      ;      \%.\par TAT\par ‘etic:\par l               c
\item \TODO{OCR 自动转写，结构图/公式需对照原 PDF 校对。}\par 4-2,      iu E       x,    7    4 o AH\par 4         m\par Sn\par 人为 Pp\par ;     <" A    =\textasciitilde{} Xic     [71\par py               也\par ran           Wk\par Mi 日   | my\par F      ay\par CaN\par \par 【第 17 页，来源：ocr】\par 六久杖目1981387888\par 上， 了 入       性\par ** 4           号 8a\par vi  we\par mm              =\par Pe.         A\par Ar   Pre |       +\par —   \%,\par WAR PPA 、  2-F* whl = - :\par veh      ae\par ly\par aM        my      —  ¥\par om\par 4      m4”\par pea\par = T ef ‘ia\par |                    i      a\par owe BT Be\par a? 2+ a\par vi     2. eye 2\par ae a    本\par aT    12                    |\par “wy”\par -的人全\par \par 【第 18 页，来源：pdf-text】
\item \TODO{OCR 自动转写，结构图/公式需对照原 PDF 校对。}\par 6-25、\par 解：依次取正对称半结构和四分之一结构如图所示，位移法基本体系见图(a)，位移法方\par 程为11\par 1\par 1\par 0\par P\par r X\par R\par +\par =\par 。画出\par 1\par M 图和M P 图如图(b)(c)所示，各系数和自由项为\par 11\par 4\par 3\par EI\par r =\par 1\par 30\par P\par R\par =\par 将系数和自由项代入位移法方程\par 1\par 4\par 30\par 0\par 3\par EI X +\par =\par ，解方程得\par 1\par 45\par 2\par X\par EI\par = −\par 。由对称性\par M=Mbar1X1 + MP叠加得原结构最后弯矩图，如图(d)所示。
\item \TODO{OCR 自动转写，结构图/公式需对照原 PDF 校对。}\par 6-26、\par 解：可将原结构荷载分解为正对称和反对称两种情况，正对称情况下取半结构，位移法基\par 本体系见图(a)，位移法方程为11\par 1\par 1\par 0\par P\par r X\par R\par +\par =\par 。画出\par 1\par M 图和M P 图如图(b)(c)所\par 示，各系数和自由项为\par 11\par 23\par 12\par EI\par r =\par 1\par 1\par 2\par P\par R\par =\par 将系数和自由项代入位移法方程\par 1\par 23\par 1\par 0\par 12\par 2\par EI X +\par =\par ，解方程得\par 1\par 6\par 23\par X\par EI\par = −\par 。由对称性\par M=Mbar1X1 + MP叠加得正对称结构弯矩图，如图(d)所示。\par \par 【第 19 页，来源：pdf-text】\par 同理反对称情况下取半结构，位移法基本体系见图(e)，取位移法方程为11\par 1\par 1\par 0\par P\par r X\par R\par +\par =\par 画出\par 1\par M 图和M P 图如图(f)(g)所示，各系数和自由项为\par 11\par 35\par 12\par EI\par r =\par 1\par 9\par 2\par P\par R\par =\par 将系数和自由项代入位移法方程\par 1\par 35\par 9\par 0\par 12\par 2\par EI X +\par =\par ，解方程得\par 1\par 54\par 35\par X\par EI\par = −\par 。由对称性\par M=Mbar1X1 + MP叠加得正对称结构弯矩图，如图(h)所示。综合以上计算，叠加正对称和反\par 对称情况下弯矩图可得原结构弯矩图如图(i)所示。\par \par 【第 20 页，来源：ocr】\par 7.64\par 3.34 231                  12 ee      189\par --                                                 t                i 1 U.96\par -                      1.03                NX              6\par 1.03                 ax                     1.2                0.86\par 6                     2.3147  No\par 3.34\par 0.519              0.514                      0.6,               0.43\par 反对称M图                        M AY (KN-m)\par (h)                                           (i)\par 41. 7/\par ON 4,\par SoZ,    ae\par oy bac e8\par OT fi \textbackslash{}O\par A a2 CRED\par BERS Asia\par \par 【第 21 页，来源：ocr】\par HEL SAR STN\par > @         )     ee\par ¥ er                   3\par fyb\par A\par i.\par Kt Paps     bg Os 4a* seh\par SFpl/18               SFpl/18\par es\par 5Fpl/l8    SFpl/18\par Bs aw) 71\par Ns\par 2F ol ag 2F el                   |\par 人\par BSi-:ReE kei\par \par 【第 22 页，来源：ocr】\par SEY SBAR Sil\par 未久徽信1981587888\par br.\par “eh TRS\par my          AAS\par (29.     ot he     th  逢>    >如\par | |  wt te Maren\par a eo\par Fle  \%,\par Me:\par Oe\par \par 【第 23 页，来源：ocr】\par IRAN o 310371665\par ho,        9   7 OX,                  3\par vas ae a\par bh\par 4\par R           ed = Fyxl= srl\par 1       As          ,\par Zl             aA hye Th hi     nh\par A\par 3H                        3\par que 1  ew 4 be  , Whe YAR.\par i\par bu (ui   « ie Ty Al 45 <-  ie\par 人py十PP        ahd\par I      ie\par he ht on =  +\par 3    :\par J        i\par | 4    :             i\par 由              j bee\par 4Juo                     ay Re  4yx jt wm       |\par nt                                                 y\par al i         Nye Bi’ im erg?  240\par oui  和\par “CRO\par BER kee\par \par 【第 24 页，来源：ocr】\par RAR Stil\par RO TUENOD 1GBUSS8\par ab                 4b\par i  :         bg | 和\par \&                                 Gy:\par OO\par \par 【第 25 页，来源：pdf-text】
\item \TODO{OCR 自动转写，结构图/公式需对照原 PDF 校对。}\par 6-32、\par 解：可将原结构支座位移分解为正对称和反对称两种情况，其中正对称下结构没有内力可不\par 考虑，取反对称半结构，由于竖杆为剪力静定杆，故水平线位移可不作为位移法基本未知量，\par 取位移法基本体系见图(a)，位移法方程为11\par 1\par 1\par 0\par P\par r X\par R\par +\par =\par 。画出\par 1\par M 图和M P 图如图(b)\par (c)所示，各系数和自由项为\par 11\par 13\par r\par i\par =\par 1\par 12\par P\par i\par R\par l\par [符号]\par = −\par 将系数和自由项代入位移法方程\par 1\par 12\par 13\par 0\par i\par iX\par l\par [符号]\par −\par =\par ，解方程得\par 1\par 12\par 13\par X\par l\par [符号]\par =\par 。由对称性\par M=Mbar1X1 + MP叠加得原结构最后弯矩图，如图(d)所示。\par \par 【第 26 页，来源：pdf-text】
\item \TODO{OCR 自动转写，结构图/公式需对照原 PDF 校对。}\par 6-32\par 解：可将支座位移分解为正对称加反对称两种情况，其中正对称下没有弯矩可不\par 考虑，取反对称半结构，位移法基本体系如图(a)所示。其中竖杆为剪力静定\par 杆，可不考虑线位移为基本未知量。画出\par 1\par M 图和\par C\par M 图，如图(b)(c)所示。\par 位移法方程为11\par 1\par 1C\par 0\par r X\par R\par +\par =\par 。各系数和自由项为：11\par 13\par r\par i\par =\par 1C\par 12i\par R\par l\par [符号]\par =\par 。将系数\par 和自由项代入位移法方程，得\par 1\par 12\par 13\par 0\par i\par iX\par l\par [符号]\par +\par =\par ，解方程得\par 1\par 12\par 13\par X\par l\par [符号]\par =\par 。利用叠加公\par 式M=\par 1\par 1\par C\par M\par M X +\par ，M 图如图(d)所示。
\item \TODO{OCR 自动转写，结构图/公式需对照原 PDF 校对。}\par 6-33\par 解：结构有3 个独立位移，加上三个附加约束，画出基本体系[见图(a)]。画出\par 1\par M 、\par 2\par M 、\par 3\par M 图和\par P\par M 图[见图(b)\textasciitilde{}(e)]，列出位移法方程\par [符号][符号]\par [符号][符号]\par [符号]\par =\par +\par +\par +\par =\par +\par +\par +\par =\par +\par +\par +\par 0\par 0\par 0\par 3\par 3\par 33\par 2\par 32\par 1\par 31\par 2\par 3\par 23\par 2\par 22\par 1\par 21\par 1\par 3\par 13\par 2\par 12\par 1\par 11\par P\par P\par P\par R\par X\par r\par X\par r\par X\par r\par R\par X\par r\par X\par r\par X\par r\par R\par X\par r\par X\par r\par X\par r\par 计算刚度系数和自由项，即\par l\par EI\par r\par 8\par 11 =\par 2\par 21\par 12\par 5\par l\par EI\par r\par r\par =\par =\par 2\par 31\par 13\par 6 l\par EI\par r\par r\par =\par =\par 3\par 22\par 13 l\par EI\par r\par =\par 0\par 32\par 23\par =\par =r\par r\par 3\par 33\par 18 l\par EI\par r\par =\par 8\par 3\par 2\par 1\par ql\par R P =\par 4\par 5\par 2\par ql\par R P\par =\par 8\par 5\par 3\par ql\par R P\par =\par \par 【第 27 页，来源：pdf-text】\par (a) 基本体系\par (b) [缺字]1图\par (c) [缺字]2图\par (d) [缺字]3图\par (e) [缺字][缺字]图\par 代入位移法方程，即\par [符号]\par [符号]\par [符号]\par [符号]\par [符号]\par [符号]\par [符号]\par [符号]\par [符号]\par =\par +\par =\par +\par =\par +\par 0\par 8\par 5\par 18\par 6\par 0\par 4\par 5\par 13\par 5\par 0\par 8\par 3\par 6\par 5\par 8\par 3\par 3\par 1\par 2\par 2\par 3\par 1\par 2\par 2\par 3\par 2\par 2\par 2\par 1\par ql\par X\par l\par EI\par X\par l\par EI\par ql\par X\par l\par EI\par X\par l\par EI\par ql\par X\par l\par EI\par X\par l\par EI\par X\par l\par EI\par [符号]\par [符号]\par [符号]\par [符号]\par [符号]\par [符号]\par [符号]\par [符号]\par [符号]\par =\par =\par =\par EI\par ql\par X\par EI\par ql\par X\par EI\par ql\par X\par 7632\par 461\par 159\par 20\par 636\par 49\par 4\par 3\par 4\par 2\par 3\par 1\par \par 【第 28 页，来源：ocr】\par EEA RS TID)\par 未久极信1981387888
\item \TODO{OCR 自动转写，结构图/公式需对照原 PDF 校对。}\par 5-34,    ay\par td        a\par mM   Si\par aa | *, gale \_ bie 有 po\par SY le     ar   ae\par 视 le a\par “a  ze Eg      人> ie\par J                oe eae\par L, HES Ke Ga et ©\par 3A\par \textasciitilde{}   a\par nee
\end{enumerate}
\chapter{第13讲：弯矩分配法}
\begin{tcolorbox}[title=解析总览]
弯矩分配法按固端弯矩、分配系数、分配和传递逐轮计算。
\end{tcolorbox}
\section{作业及答案：结构力学基础与技巧班第13次作业及答案（弯矩分配法）}
\begin{enumerate}[label=\textbf{\arabic*.},leftmargin=2em]
\item \TODO{OCR 自动转写，结构图/公式需对照原 PDF 校对。}\par 【第 1 页，来源：ocr】\par 汞久向信1981287888\par 七、练习题\par 1， 基本概念
\item \TODO{OCR 自动转写，结构图/公式需对照原 PDF 校对。}\par 7-1 图7-51中哪一种情况不能用力矩分配法计算。(  ) (河海大学 2005)\par |     A          |     B          上. C         |     D\par 图7-51
\item \TODO{OCR 自动转写，结构图/公式需对照原 PDF 校对。}\par 7-2 7- 52 所示结构AC杆A 端的力矩分配系数puac王\_\_\_ 。(浙江大学 2009)
\item \TODO{OCR 自动转写，结构图/公式需对照原 PDF 校对。}\par 7-3 7-53 所示结构，要使结点 中产生单位转角，则在结点 妃需施加的外力偶为\par ( ， )。(合肥工业大学 2010、西南交通大学 2010)\par A. 13;                B. 10i                 C. 9i                  D. 8i
\item \TODO{OCR 自动转写，结构图/公式需对照原 PDF 校对。}\par 7-4 7-54 所示刚架各杆的线刚度 ;一常数，要使 A 点发生顺时针的单位转角，则\par A 点需要施加的弯矩mm一。 (广西大学 2010)\par uot 1 '; 4\par ON ‘N : 420.\par 254, ; o ASS\par ONO\par 萌锡线上快印\par \par 【第 2 页，来源：ocr】\par 后绢人太基公众引胡与|\par 东久币信19813587888\par C\par 2.5EI               g\par A                    D                         m\par 2\&1                                                                    A\par EI                   §          4\textbackslash{}Gi             3i ML ©                                    |\par B                                              B\par |      4m      |                 :              Jo |      /|\par 图7-52                  图7-53                    图7-54\par 2. 荷载作用下的计算
\item \TODO{OCR 自动转写，结构图/公式需对照原 PDF 校对。}\par 7-5 图7-55 maw, SHEL 为常数，截面C、D 两处的弯矩值 Mc、Mp 分别为\par (    )。(单位: kN。m) 〈哈尔滨工业大学 2009)\par A. 1.0，2.0       B. 2.0, 1.0       Cc. 一1.0，一2.0 D. —2.0, —1.0
\item \TODO{OCR 自动转写，结构图/公式需对照原 PDF 校对。}\par 7-6 用力矩分配法求图 7- 56 Pra, i=1. (MK 2012)
\item \TODO{OCR 自动转写，结构图/公式需对照原 PDF 校对。}\par 7-7 用力矩分配法计算图 7 - 57 所示结构，并作 M 图。 (计算两轮) 〈哈尔滨工业大\par 学 2009)\par C| 9mN             be                                      40KN/m\par 5                                        LITT TT ttt\}\par \&             vy                              1261\par P           B                              El        1.2E1        1,      §\par D                    A                                                    ，\par |  3m    |  ,'™                |     1     |                |   3m    |       6m        |\par 图7-55             图7-56                  图7-57
\item \TODO{OCR 自动转写，结构图/公式需对照原 PDF 校对。}\par 7-8 计算图7- 58 所示结构，并绘 M 图。要求在力法、位移法、力矩分配法中选择，\par EI 为常数。(重庆大学 2010)
\item \TODO{OCR 自动转写，结构图/公式需对照原 PDF 校对。}\par 7-9 用力矩分配法计算图 7 - 59 所示结构，并作 M 图。已知: g=2.4kN/m, \&4F\par ET相同。(每结点分配两次) (西南交通大学 2008)\par |"  Cc                               q\par A\par 287            El         §                   B               c            r\par A           B                                                                  F\par EI              \&\par E              D         G\par |    8m     |  4m |  4m  |                 |   7.5m   |      10m      |   7.5m   |\par 图7-58 7-59\par Ne 7 ' iff\par O75 \textbackslash{}O\par 萌锡线上快印\par \par 【第 3 页，来源：ocr】\par EAR Stil\par RATE OD ICHIBSS
\item \TODO{OCR 自动转写，结构图/公式需对照原 PDF 校对。}\par 7-10 用力矩分配法作图 7 - 60 FRERBRHSEA. REI=KHR. (山东建筑大学\par 2012)
\item \TODO{OCR 自动转写，结构图/公式需对照原 PDF 校对。}\par 7-11 用力矩分配法求解图 7 - 61 所示结构，                 eg             10KN\par 作 M图。FET 为常数〈忽略前力和轴力的影响)。 F    A       B    C     D
\item \TODO{OCR 自动转写，结构图/公式需对照原 PDF 校对。}\par 7-12 用力矩分配法计算图 7 - 62 所示结构，            图7-60\par 并作 M图，各杆 ETI 相同，为常数。(重庆大学 2011)\par 2kN/m            3kN/m                            N/m\par 6kN-m\par 4    El    C ET    D El F\par |         B        c      E   \&\par 8kN\par B                      A        D°         \&\par |   4m    |  3m  |    4m            |   4m    |    4m    |\par 图7-61                                  图7-62\par 八、练习题答案
\item \TODO{OCR 自动转写，结构图/公式需对照原 PDF 校对。}\par 7-1 D。
\item \TODO{OCR 自动转写，结构图/公式需对照原 PDF 校对。}\par 7-2 4/7。提示: C支座相当于固定支座。
\item \TODO{OCR 自动转写，结构图/公式需对照原 PDF 校对。}\par 7-3 B。提示: BC 杆转动刚度为零。
\item \TODO{OCR 自动转写，结构图/公式需对照原 PDF 校对。}\par 7-4 1lz。
\item \TODO{OCR 自动转写，结构图/公式需对照原 PDF 校对。}\par 7-5 B。
\item \TODO{OCR 自动转写，结构图/公式需对照原 PDF 校对。}\par 7-6 M图见图7-63。
\item \TODO{OCR 自动转写，结构图/公式需对照原 PDF 校对。}\par 7-7 M图见图7- 64。提示: Mi, = X3? =45KN 。m。\par “     115.65\par ge                  \$8.59                81.14\par ossee                            47.1    B\par 0.29? [06soe\par A    .\par O.1gP                          23.55         40.57\par MA                            NM图(单位: KN-m)
\item \TODO{OCR 自动转写，结构图/公式需对照原 PDF 校对。}\par 7-63 7-64
\item \TODO{OCR 自动转写，结构图/公式需对照原 PDF 校对。}\par 7-8 M图见图7- 65。提示: 本题用力矩分配法较方便，C 支座相当于固定支座。
\item \TODO{OCR 自动转写，结构图/公式需对照原 PDF 校对。}\par 7-9 M图见图7- 66。提示: CFG 为静定部分，该部分弯和矩图可以直接画出，并且其\par 在 C 点的反作用力对左侧超静定部分的弯矩图无影响。\par wha.\par Ne 。  A\par C7 i VO\par fA eee FIRED\par \par 【第 4 页，来源：ocr】\par SEA Rains]\par RPAEVUENPD ABUSOS\par 49\par 17.36\par 8.68               12.08\par 12.08\par 0 >                                          8.68\par 45                                                          16.88\par 15\par 434           6.04\par M¥A(xFp/41)                                                                                                                           2M图(单位: KN-m)
\item \TODO{OCR 自动转写，结构图/公式需对照原 PDF 校对。}\par 7-65 7-66
\item \TODO{OCR 自动转写，结构图/公式需对照原 PDF 校对。}\par 7-10 MARA 7-67.
\item \TODO{OCR 自动转写，结构图/公式需对照原 PDF 校对。}\par 7-11 M图见图7-68。\par 12 857      i\par 4                    3.43\par 25\par AM图(单位: KN-m)                                                                                     AM图(单位: KN-m)
\item \TODO{OCR 自动转写，结构图/公式需对照原 PDF 校对。}\par 7-67 7-68
\item \TODO{OCR 自动转写，结构图/公式需对照原 PDF 校对。}\par 7-12 MBL 7-69.\par 16\par 27/8\par 5\par 4           6\par 2图(单位: Nm)
\item \TODO{OCR 自动转写，结构图/公式需对照原 PDF 校对。}\par 7-69\par voltae\par 。  、\textbackslash{} < 《LA 一\par ra 77.  ww e   \textasciitilde{}\par ? /; ' 1 A)
\end{enumerate}
\chapter{第14讲：矩阵位移法}
\begin{tcolorbox}[title=解析总览]
矩阵位移法按单元刚度、坐标转换、总体组装、边界条件和回代内力求解。
\end{tcolorbox}
\section{作业及答案：结构力学基础与技巧班第14次作业及答案（矩阵位移法）}
\begin{enumerate}[label=\textbf{\arabic*.},leftmargin=2em]
\item \TODO{OCR 自动转写，结构图/公式需对照原 PDF 校对。}\par 【第 1 页，来源：ocr】\par ESA Raia]\par ROE VEN OT CBISCS\par 八、练习题\par 1，基本概念
\item \TODO{OCR 自动转写，结构图/公式需对照原 PDF 校对。}\par 8-1 单元刚度矩阵是指下列两组量值之间关系的变换矩阵。( ) (北京科技大学\par 2010、中国地震局工程力学研究所 2009、中南大学 2010)\par A. 杆端位移和杆端力             B. 杆端位移和结点位移，\par C. 杆端位移和结点荷载           D. 结点位移和结点荷载
\item \TODO{OCR 自动转写，结构图/公式需对照原 PDF 校对。}\par 8-2 判断题。结构刚度矩阵反映了结构结点位移与荷载之间的关系 。( ) (烟台\par 大学 2013)\par 2. 刚度矩阵的计算
\item \TODO{OCR 自动转写，结构图/公式需对照原 PDF 校对。}\par 8-3 试用先处理法建立图 8 - 43 所示结构的结构刚度矩阵，忽略杆件的轴向变形\par (TI一常数) 。(湖南大学 2008)
\item \TODO{OCR 自动转写，结构图/公式需对照原 PDF 校对。}\par 8-4 求图8- 44 所示刚架结构刚度矩阵的任意两个主元素〈自由结点)。已知各杆\par IJ/A=1/10。(西南交通大学 2010)\par 5          8\par 2        4       5            4         7\par @         @   |        ;         >\par (0)       ®       是\par 图8-43                     图8-44\par 3. 等效结点荷载
\item \TODO{OCR 自动转写，结构图/公式需对照原 PDF 校对。}\par 8-5 图8-45 所示结构按和矩阵位移法计算，则与结点位移 1、2 〈正方向见图虚线标\par 示) 对应的等效结点荷载向量为\_\_\_\_\_“。(浙江大学 2012)
\item \TODO{OCR 自动转写，结构图/公式需对照原 PDF 校对。}\par 8-6 8-45 所示结构，已知两杆 已I一常数。若按矩阵位移法求得结点位移 1、2 分\par 别为和6、A:，则杆件 BC 左、右端的弯矩各为 、\_ (均以下侧受拉为正)。\par (浙江大学 2012)
\item \TODO{OCR 自动转写，结构图/公式需对照原 PDF 校对。}\par 8-7 试求图 8 - 46 所示结构中结点 3 的综合结点荷载列阵 F,= (X， Y; Ms)".\par (中南大学 2009)\par wha.\par Ne 。  A\par O75 \textbackslash{}O\par 萌锡线上快印\par \par 【第 2 页，来源：ocr】\par ES RA Rall\par RAEVENODIGSISSS\par 12KN/m\_\_iskN-m | 20kN\par 12             2             3          4            5                   ,\par qd \_-s!       q       1                               7                        P\par In   , 2        1      4m     4m     4m\par FA 8-45    fl 8 - 46
\item \TODO{OCR 自动转写，结构图/公式需对照原 PDF 校对。}\par 8-8 按先处理法求图 8 - 47 所示连续梁的结点荷载列阵 P。(西南交通大学 2007)
\item \TODO{OCR 自动转写，结构图/公式需对照原 PDF 校对。}\par 8-9 图8-48 (a) 所示结构，不考虑\par 轴向变形，整体坐标见图 8 - 48 (b)，图中                         20kNm\par 括号内数码为结点总码 (CHAMBRE, sev 08) oO)          uN\par 竖直、转动方向顺序排列) ，求等效结点荷载\par PURE P。(山东建筑大学 2008)\par 2kN         4kN                           skim                         §                .\par km        12kN/m\par Mm,\par ?可      al           i           >  ,                       4m\par x             |一 |\par |                                 @                         0)
\item \TODO{OCR 自动转写，结构图/公式需对照原 PDF 校对。}\par 8-47                                                                    图8-48
\item \TODO{OCR 自动转写，结构图/公式需对照原 PDF 校对。}\par 8-10 8-49 所示平面结构用和抢阵位移法计算，引入支承条件后的整体刚度矩阵为\par 多少阶? 求结点 2 和结点 5 的综合结点荷载列阵。(中南大学 2011)\par 4, 内力和位移的计算\par :      -          8-11 已知图 8-50 所示连续梁的结点\par @                   位移列阵 4一(一0.003 0.004 0.005)T (分\par \textasciitilde{}          >? .°             别为结点 2、3、4 的转角)，设各杆的 下一2X\par ake|2     ? 5           104kNVcm2: ，T一400cm4 ，!一400cm。求单元\par 二 @                  ORM AE, HRT AA 2010)\par ©\par \textasciitilde{}          0)\par 了 1                            7    2.4kKN/m    4kN-m     5.6kN-m\par .        Ree,                1 0    二      @    四         ®   和
\item \TODO{OCR 自动转写，结构图/公式需对照原 PDF 校对。}\par 8-49 8-50
\item \TODO{OCR 自动转写，结构图/公式需对照原 PDF 校对。}\par 8-12 ARBOR 8-51 所示检架各杆轴力。已知: EA=60kN, Fp=100kN\par (上海大学 2008、武汉大学 2008 类似)
\item \TODO{OCR 自动转写，结构图/公式需对照原 PDF 校对。}\par 8-13 8-52 所示结构，已知结点 2 的水平、竖向和转角位移分别为 0、一0. 0926、\par 一0. 1667，结点 3 的转角为 0.2220，试求 23 FF 2 端杆端力。已知单元刚度矩阵中的部分\par wha.\par en 2,\par C7 i A)\par BA Sete FIRED\par \par 【第 3 页，来源：ocr】\par ES EAR Sips)\par ROE VEN OT CBISCS\par 元素:\par ®      Fp\par 1kN\par Ig=2m*          im‘          上 M,@\par 4m                    lm       lm          *\par \$——"           全一 th 一\par 图8-51                     图8-52\par 学 2008)
\item \TODO{OCR 自动转写，结构图/公式需对照原 PDF 校对。}\par 8-14 用和矩阵位移法计算并作图 8 - 53 所示连续梁的弯矩图。各杆 下IT 相同，为常数。\par (重庆大学 2009)\par 20kN       10kN/m    30kN-m       10kN\par 3m  \}  3m  '     6m.    +    6m    \{\par FA 8-53\par 九、练习题答案
\item \TODO{OCR 自动转写，结构图/公式需对照原 PDF 校对。}\par 8-1 A.
\item \TODO{OCR 自动转写，结构图/公式需对照原 PDF 校对。}\par 8-2 X, Bm: 反映的是结点位移和结点荷载之间的关系。\par EI/18 EI/6 0
\item \TODO{OCR 自动转写，结构图/公式需对照原 PDF 校对。}\par 8-3 K=| EI/6 4EI/3 ETI/3|。提示: 建立顺时针坐标系。\par 0 FEI/3 2EI
\item \TODO{OCR 自动转写，结构图/公式需对照原 PDF 校对。}\par 8-4 Ky=43EI/16, Kss=43EI/16, Ke=2EI. 2€: (1) AMAA HRA,\par ae hae Mat hE, BLA EMS —Ae MR, BOLT a oh eH a\par HE, (2) 坚杆单元需要坐标变换，若假设单元方向向上，则 a一90"。
\item \TODO{OCR 自动转写，结构图/公式需对照原 PDF 校对。}\par 8-5 一q12/24; ql/2.\par 4EI0, 6EIA, gl? \_ 2EI@, , 6EIA; ql?
\item \TODO{OCR 自动转写，结构图/公式需对照原 PDF 校对。}\par 8-6 FG    a 453 — te Tee
\item \TODO{OCR 自动转写，结构图/公式需对照原 PDF 校对。}\par 8-7 F;= (0 —34kN 一9kN。m)7 (假设杆端弯矩逆时针为正) 。
\item \TODO{OCR 自动转写，结构图/公式需对照原 PDF 校对。}\par 8-8 P=P;+Pp= (0 0 一16)7十 (一2 一5 0)T= (一2kN 一5kN 一16kN 。m)T。\par 注意: 题中 4kN 的结点力作用在支座处，对结点荷载列阵无影响。
\item \TODO{OCR 自动转写，结构图/公式需对照原 PDF 校对。}\par 8-9 P=Pp+Pp= (12 一8 0)T+ (一5十3 0 20)7= (10kKN —8kN+m\par 20kN 。m)TI。\par vat 1 ik\par ote\par O75 \textbackslash{}O\par BA Sete FIRED\par \par 【第 4 页，来源：ocr】\par ES EAR Stil\par RATE OBICDUSSS
\item \TODO{OCR 自动转写，结构图/公式需对照原 PDF 校对。}\par 8-10 12R. Pea= (ql —ql/2 0)7, Ps= (0 —ql/2 ql*/8)? (假设杆端弯矩逆\par 时针为正) 。\par 5. 1kN
\item \TODO{OCR 自动转写，结构图/公式需对照原 PDF 校对。}\par 8-11  -| oe 提示: 原题未规定结点转角的正方向，读者可以根据荷\par 一1. 2kN 。m\par 载情况判断出 4 结点转角应该为道时针，而已知条件中 4 结点的转角为 0.005 〈正值)，由\par JETER EAE. BA 8 - 54 所示坐标系。
\item \TODO{OCR 自动转写，结构图/公式需对照原 PDF 校对。}\par 8-12 各杆轴力见图8- 55。
\item \TODO{OCR 自动转写，结构图/公式需对照原 PDF 校对。}\par 8-13 (0 一0.7794E ”一0.7784E) 。提示: 用公式下=大4 。题中未给出单元刚度\par 矩阵中的第一行元素，需要读者自行记住。
\item \TODO{OCR 自动转写，结构图/公式需对照原 PDF 校对。}\par 8-14 M图见图8- 56。\par 70.37       \_\par ° -22.22                  23             2\par ll\par M,@\par ae           RN(单位: KN)                      J.          M图(单位: KN-m)\par 图8-54        图8-55                      图8-56\par Vy.\par \textasciitilde{}. 1g 2. aA\par O75 \textbackslash{}O\par BA Sete FIRED
\end{enumerate}
\backmatter
\chapter*{OCR 处理报告}
\begin{itemize}
\item 第1讲 结构力学基础与技巧班第1次作业答案（几何组成分析）.pdf: 6页, \{'ocr': 0, 'pdf-text': 6, 'empty': 0\}
\item 第1讲 结构力学基础与技巧班第1次作业（几何组成分析）.pdf: 3页, \{'ocr': 3, 'pdf-text': 0, 'empty': 0\}
\item 第2讲 结构力学基础与技巧班第2次作业答案（理论力学回顾）.pdf: 4页, \{'ocr': 0, 'pdf-text': 4, 'empty': 0\}
\item 第2讲 结构力学基础与技巧班第2次作业（理论力学回顾）.pdf: 2页, \{'ocr': 2, 'pdf-text': 0, 'empty': 0\}
\item 第3讲 结构力学基础与技巧班第3次作业答案（材料力学回顾）.pdf: 7页, \{'ocr': 0, 'pdf-text': 7, 'empty': 0\}
\item 第3讲 结构力学基础与技巧班第3次作业（材料力学回顾）.pdf: 12页, \{'ocr': 12, 'pdf-text': 0, 'empty': 0\}
\item 第4讲 结构力学基础与技巧班第4次作业及答案（静定梁和刚架）.pdf: 14页, \{'ocr': 14, 'pdf-text': 0, 'empty': 0\}
\item 第5讲 结构力学基础与技巧班第5次作业及答案（静定桁架）.pdf: 3页, \{'ocr': 3, 'pdf-text': 0, 'empty': 0\}
\item 第6讲 结构力学基础与技巧班第6次作业及答案（组合结构）.pdf: 3页, \{'ocr': 3, 'pdf-text': 0, 'empty': 0\}
\item 第7讲 结构力学基础与技巧班第7次作业及答案（三铰拱）.pdf: 1页, \{'ocr': 1, 'pdf-text': 0, 'empty': 0\}
\item 第8讲 结构力学基础与技巧班第8次作业及答案（静定结构影响线）.pdf: 18页, \{'ocr': 14, 'pdf-text': 4, 'empty': 0\}
\item 第9讲 结构力学基础与技巧班第9次作业及答案（静定结构位移计算）.pdf: 4页, \{'ocr': 4, 'pdf-text': 0, 'empty': 0\}
\item 第10讲 结构力学基础与技巧班第10次作业及答案（力法1）.pdf: 21页, \{'ocr': 20, 'pdf-text': 1, 'empty': 0\}
\item 第11讲 结构力学基础与技巧班第11次作业及答案（力法2）.pdf: 8页, \{'ocr': 7, 'pdf-text': 1, 'empty': 0\}
\item 第12讲 结构力学基础与技巧班第12次作业及答案（位移法）.pdf: 28页, \{'ocr': 23, 'pdf-text': 5, 'empty': 0\}
\item 第13讲 结构力学基础与技巧班第13次作业及答案（弯矩分配法）.pdf: 4页, \{'ocr': 4, 'pdf-text': 0, 'empty': 0\}
\item 第14讲 结构力学基础与技巧班第14次作业及答案（矩阵位移法）.pdf: 4页, \{'ocr': 4, 'pdf-text': 0, 'empty': 0\}
\end{itemize}
\end{document}
